\documentclass[author={Karen Coutinho, Ishaan Bhalla, Aman Raj, Dixit Kaloorani Malarmannan},
                supervisor={Dr. Daniel D'Andrea},
                degree={MSc},
                title={CellSelector Omics: A Multi-Omics Recommendation System for Cell Line Selection},
                subtitle={},
                type={},
                year={2026}]{dissertation}

\begin{document}

\let\cleardoublepage\clearpage

\maketitle

\frontmatter

\makedecl
\makedeclai

\tableofcontents
\listoffigures
\listoftables

% =============================================================================

\chapter*{Abstract}

Cell line selection in cancer research requires manual cross-referencing of incompatible databases with conflicting identifiers, units and sparse coverage. CellSelector Omics solves this by integrating 2,076 cell lines and 222 million expression records across five sources (HPA, DepMap, GEO, CCLE proteomics, KEGG), including partially-covered lines.

Key finding: gene class determines ranking accuracy more than global weighting schemes. Tissue-specific genes (EGFR, ERBB2) rank accurately by expression level. Ubiquitous genes (GAPDH, ACTB) cannot be ranked by expression alone. Loss-of-function genes (TP53, BRCA1) rank by mutation status. Per-class weighting improves Mean Reciprocal Rank from 0.1096 to 0.2367 (+116\%, p=0.0006). Loss-of-function genes achieve 4,700\% improvement (0.0041 → 0.1968). Precision@k confirms performance: tissue-specific P@1=0.75, P@5=0.63; loss-of-function P@1=0.44, P@5=0.34.

Without biological constraints, learned weights converge to biological principles. Loss-of-function genes receive mutation-impact weight 0.80 and expression weight 0.003 (270-fold difference). This is exactly what experts anticipate. The convergence is not accidental; rather, it indicates that gene class is a fundamental property of the ranking problem that may be found.

Three severe problems were found through systematic validation: score-saturation disguising weak signal, non-deterministic tie-breaking violating reproducibility, and misattribution of HPA/GEO identifiers (94\% and 80\% inaccurate). A biological flaw was identified by domain-expert validation at AstraZeneca: wild-type TP53-positive lines performed worse than TP53 mutants. A TP53-specific penalty corrected this, improving ubiquitous-class MRR from 0.0633 to 0.1216. This outcome demonstrates that face-validity assessment against domain expertise is necessary alongside statistical evaluation in biomedical tool development.

Pathway-coherence analysis tested four methods (naive co-expression, ssGSEA union/maximum, specificity-weighted combination); all failed to improve ranking. Mechanistic analysis shows phosphorylation cascades do not produce coordinated transcriptional upregulation and ssGSEA methods cannot detect phosphorylation-based activation. Understanding this failure motivated mutation-impact scoring.

System: FastAPI backend and React frontend with per-source evidence visualisation, both deployed on a single Google Cloud Platform Compute Engine VM (e2-standard-2, asia-south2-b). Backend runs via Docker container managed by docker-compose; frontend built into static files and served by Nginx on the same VM. Neo4j knowledge graph stored in a separate self-hosted Docker container; expression and proteomics data in Parquet files on local disk. Deterministic ranking (seed=42, stable sort) for bitwise reproducibility. Manual deployment via SSH (git pull, docker-compose rebuild, npm build and nginx reload). All recommendations include evidence-backed citations (PubMed, KEGG, Cellosaurus, omics data).

Results validated on independent AstraZeneca test genes. Limitations: no experimental validation (priority future work), validation set n=44, proteomics coverage 18.1\%, mutation annotations 60.1\%, cancer cell lines only. Future work includes experimental validation, OncoKB/COSMIC expansion, tissue-specific reliability assessment and extension to non-cancer contexts.

When deciding which data modality should rank each gene, gene class emerges as a crucial criteria. The scoring strategy is determined by biological function rather than global optimization. This insight directly informs multi-omics integration: single weighting schemes are insufficient; the ranking method must match functional biology.
% =============================================================================

\chapter*{Supporting Technologies}

\textbf{Backend and API.}
\begin{itemize}
    \item FastAPI to build the backend API, with Uvicorn as the ASGI server and Pydantic to validate request and response payloads.
\end{itemize}

\textbf{Data processing and ranking.}
\begin{itemize}
    \item PyArrow to write harmonised data into Parquet files with Snappy compression.
    \item pandas and NumPy to read Parquet files directly at runtime, percentile and rank-fusion calculations used throughout the ranking algorithm.
    \item SciPy (\texttt{scipy.optimize}) to learn the per-gene-class evidence weights via constrained numerical optimisation against Mean Reciprocal Rank.
    \item NetworkX to implement Random Walk with Restart via personalised PageRank over the knowledge graph.
\end{itemize}

\textbf{Knowledge graph.}
\begin{itemize}
    \item Neo4j to store and query the multi-hop knowledge graph linking genes, cell lines, pathways and mutations, deployed in its own Docker container on the same VM.
\end{itemize}

\textbf{Agentic justification pipeline.}
\begin{itemize}
    \item The Groq API, serving the \texttt{openai/gpt-oss-20b} model, for retrieval-augmented generation of citation-backed justifications, with a local Ollama fallback (\texttt{llama3.1:8b}, falling back further to \texttt{llama3:latest}) used when the hosted endpoint is unavailable.
    \item The PubMed E-utilities and KEGG REST APIs for literature and pathway retrieval feeding the justification pipeline.
\end{itemize}

\textbf{Frontend.}
\begin{itemize}
    \item React, TypeScript and the Vite build tool to build the single-page frontend, styled with Tailwind CSS.
\end{itemize}

\textbf{Infrastructure and deployment.}
\begin{itemize}
    \item Google Cloud Platform Compute Engine (e2-standard-2, zone asia-south2-b) to provide a single virtual machine hosting all application components.
    \item Docker and docker-compose to containerise the FastAPI backend and Neo4j knowledge graph, both deployed and running on the same VM.
    \item The frontend built as static files (dist/) via npm run build and served directly by Nginx on the VM.
    \item Nginx to reverse-proxy API requests to the backend Uvicorn service and serve the compiled frontend.
    \item Certbot with Let's Encrypt to obtain and renew HTTPS certificates.
    \item No automated CI/CD pipeline. Deployment is entirely manual, performed over SSH by executing git pull on the VM, docker-compose up -d --build to rebuild and restart the backend and Neo4j containers, and npm run build followed by nginx reload to update the frontend. Typical deployment time is a few minutes depending on the magnitude of changes.
\end{itemize}

\textbf{Version control and deployment.}
\begin{itemize}
    \item Git and GitHub for version control. Deployment is entirely manual, executed over SSH by running git pull on the VM, then docker-compose up -d --build to rebuild and restart the backend and Neo4j containers, followed by npm run build and nginx reload to update the frontend.
\end{itemize}

\textbf{Data sources and storage.}
\begin{itemize}
    \item Five publicly available datasets integrated end-to-end: the Human Protein Atlas, DepMap (expression and CRISPR gene dependency), the Gene Expression Omnibus, CCLE Proteomics and KEGG Pathways, harmonised against canonical Cellosaurus and Ensembl identifiers as described in Section~2.3.
    \item Parquet files containing harmonised expression, proteomics and derived ranking data stored directly on the VM's local disk.
\end{itemize}

\textbf{Documentation.}
\begin{itemize}
    \item LaTeX to format this dissertation, via the online service Overleaf.
\end{itemize}
% =============================================================================
\chapter*{Notation and Acronyms}
\addcontentsline{toc}{chapter}{Notation and Acronyms}

{\footnotesize
\renewcommand{\arraystretch}{1.1}
\setlength{\tabcolsep}{6pt}

\begin{tabular}{@{}p{2.3cm}cp{11.5cm}@{}}
ACH      & : & DepMap cell line accession prefix, for example ACH-000164 \\
API      & : & Application Programming Interface \\
CCLE     & : & Cancer Cell Line Encyclopedia \\
CI/CD    & : & Continuous Integration / Continuous Deployment \\
CRISPR   & : & Clustered Regularly Interspaced Short Palindromic Repeats \\
CSV      & : & Comma-Separated Values \\
CVCL     & : & Cellosaurus cell line accession prefix, for example CVCL\_0063 \\
DepMap   & : & Cancer Dependency Map \\
ENSG     & : & Ensembl stable gene identifier prefix, for example ENSG00000005471 \\
GEO      & : & Gene Expression Omnibus \\
HPA      & : & Human Protein Atlas \\
HTTPS    & : & Hypertext Transfer Protocol Secure \\
JSON     & : & JavaScript Object Notation \\
KEGG     & : & Kyoto Encyclopedia of Genes and Genomes \\
LLM      & : & Large Language Model \\
LOO-CV   & : & Leave-One-Out Cross-Validation \\
miRNA    & : & MicroRNA \\
MRR      & : & Mean Reciprocal Rank \\
MSI      & : & Microsatellite Instability \\
P@k      & : & Precision at rank $k$ \\
PDX      & : & Patient-Derived Xenograft \\
RAG      & : & Retrieval-Augmented Generation \\
RGCN     & : & Relational Graph Convolutional Network \\
RNA      & : & Ribonucleic Acid \\
RRF      & : & Reciprocal Rank Fusion \\
RWR      & : & Random Walk with Restart \\
SCLC     & : & Small Cell Lung Cancer \\
ssGSEA   & : & single-sample Gene Set Enrichment Analysis \\
TPM      & : & Transcripts Per Million \\
UI       & : & User Interface \\
\end{tabular}

\vspace{0.4cm}

\begin{tabular}{@{}p{2.6cm}cp{11.2cm}@{}}

$\mathrm{RRF}(r)$       & : & the reciprocal rank fusion score of a cell line at rank $r$ \\
$\mathrm{MRR}$          & : & Mean Reciprocal Rank, the mean of $1/r$ over each gene's correct-answer rank \\
$\mathrm{P@}k$          & : & the fraction of the top-$k$ ranked cell lines that are correct \\
$p\text{-value}$        & : & the probability, under the paired bootstrap null hypothesis, of observing an improvement at least as large as the one measured \\
\end{tabular}
}


\mainmatter

% =============================================================================

\chapter{Introduction}
\label{chap:introduction}

\section{Motivation}

Cell lines, immortalised populations of cells derived from human tissue, are the foundational experimental substrate for biopharmaceutical research. A cell line's resemblance to the biological context of interest determines how far experimental results can be trusted and generalised. Poor cell line selection invalidates assays, undermines reproducibility and wastes months of laboratory time and material cost.

In practice, cell line selection relies on manual cross-referencing of DepMap, Human Protein Atlas and Gene Expression Omnibus. Each database uses different formats, identifier schemes and units of measurement, capturing only a partially overlapping set of cell lines. This fragmentation creates friction, inconsistency and the risk that researchers will miss critical information because it exists in a database they did not think to consult.

This manual method is a real bottleneck at companies like AstraZeneca. Research scientists spend hours cross-referencing spreadsheets, manually downloading data, attempting to reconcile expression values on incompatible scales, searching for missing information. A scientist selecting a cell line for a study involving EGFR and KRAS is forced to check DepMap for expression levels and CRISPR dependency, check HPA for protein-level data and tissue specificity, check GEO for additional confirmatory expression data, cross-reference nomenclature to ensure they are looking at the same physical cell line, manually reconcile values on different scales and document which evidence came from which source. This process is error-prone, time-consuming and produces inconsistent results. The same gene searched on different days or by different scientists can yield different recommendations. A cell line that should have been selected is missed because it was named differently in one database.

\section{The CellSelector Omics Project}

CellSelector Omics addresses this problem directly. Given one or more target genes, the tool returns a ranked list of candidate human cell lines with molecularly similar alternatives. Instead of an opaque score, each recommendation is paired with an explanation supported by citations and evidence. A scientist can see not just which cell lines are recommended, but why they were chosen, what evidence is strong or thin, where the recommendation is confident and where it relies on sparse data.

The system integrates five public databases into a single queryable structure: Human Protein Atlas (HPA) with RNA-seq normalised expression across 1,198 cell lines, Cancer Dependency Map (DepMap) with expression and CRISPR gene dependency across 1,399 cell lines, Gene Expression Omnibus (GEO) with heterogeneous third-party expression data from 3,159 datasets, CCLE Proteomics with mass-spectrometry protein abundance across 375 cell lines and KEGG Pathways with biological pathway annotations linking genes to cellular processes.

This integration spans across 2,076 unique cell lines and 222 million expression records. More importantly, it does this without discarding partially-covered cell lines. A cell line present in some sources but not others is ranked fairly based on the evidence that exists for it, rather than penalised for missing data from sources that do not cover it.

The system operates through two complementary mechanisms. First, a classical statistical model scores cell lines via a transparent, interpretable ranking algorithm where each evidence type contributes a named, weighted component. Second, an autonomous AI model uses retrieval-augmented generation to produce structured justifications that cite PubMed, KEGG, Cellosaurus and the underlying omics data. This separation ensures that the ranking is reproducible and explainable, while the language model's role is purely to explain, not to decide.

\section{Central Challenges}

Three core technical challenges had to be solved to make this system work.

{\bf Challenge 1: Data Heterogeneity.} The source datasets use different naming conventions, gene identifiers and expression units. DepMap provides log-transformed TPM values. HPA provides TPM. GEO provides linear values whose original normalisation cannot be verified. The same numeric value means something different in each. Averaging across them directly produces a quantity with no interpretation. Converting between units is insufficient, since values on a shared scale remain sensitive to the sample they were drawn from. The same physical cell line is referred to by different names: HCC827, Hcc-827 and DepMap's ACH-000164 all denote the same line. Resolving every reference onto a single canonical identifier without losing any cell lines required substantial data engineering.

{\bf Challenge 2: Interpretability.} A ranking is of limited use unless a scientist can understand why a given cell line was recommended. Black-box machine learning approaches might achieve higher accuracy but would fail this requirement entirely. The system needed to expose per-source evidence so that a scientist could see not just that a cell line was recommended, but exactly why, down to the source dataset or literature reference. This necessitated stringent limitations on the agentic AI component and an open scoring pipeline.

{\bf Challenge 3: Biological Validity.} Fundamentally distinct ranking logic is needed for different kinds of genes. Because they are highly and selectively expressed in certain malignancies, receptor tyrosine kinases like EGFR are investigated in great detail. Models that are truly helpful are found by ranking cell lines according to their raw EGFR expression. Tumor suppressor genes, such as TP53, are expressed at roughly identical levels in nearly all cell lines and are examined by their absence or mutation rather than their expression. For this class, raw-expression ranking is either deceptive or uninformative. These differences cannot be captured by a single global weighting method. One of the main conclusions of the project was the early identification of this need and the methodical investigation of it.

\section{Aims, Objectives and Achievements}

{\bf Aim 1: Integrate Multiple Public Multi-Omics Datasets into a Single Queryable Structure.}

The goal is to unify multiple separate data sources using canonical cell line IDs. So that a single UI could be used to query them.

Success: Traditional Cellosaurus identifiers were mapped to 2,076 cell lines and 222 million expression entries from five databases (HPA, DepMap, GEO, CCLE, and KEGG). This integrated data warehouse allows for quick querying of the entire database on any standard computer over Parquet files using pandas.

{\bf Aim 2: Enable Gene-Presence and Gene-Exclusion Driven Ranking.}

Goal: Answering a question where a scientist specifies genes that need to be expressed (including) and those which should be expressed at all (excluding).

Achievement: There is complete support for both. Ranked cell lines arranged according to evidence of target gene expression are returned by gene-inclusion inquiries. By excluding confusing cell lines from the findings, gene-exclusion queries enable researchers to focus on models appropriate for their particular experimental setup. Treating exclusion criteria as limitations rather than preferences, exclusion functions as hard filtering as opposed to soft penalties.

{\bf Aim 3: Quantify Confidence in Each Recommendation Despite Sparse, Heterogeneous Data.}

Goal: Instead of concealing ambiguity, bring it to light. Create confidence estimates that clearly indicate areas of sparse, conflicting, or absent evidence.

Accomplishment: A gene-class aware scoring model with cross-source consistency checks and explicit handling of missing data. The approach tracks coverage independently and averages scores across various sources rather than penalizing cell lines that are not covered in certain sources. This allows scientists to see the quality of the supporting evidence as well as the ranking. Per-gene-class weighting increases ranking accuracy by 116\% (from MRR 0.1096 to 0.2367), with the biggest gains for loss-of-function genes where the key signal is mutation status rather than expression.

{\bf Aim 4: Provide Transparent, Evidence-Cited Justifications.}

Goal: Each suggestion should be accompanied by a well-organized, citation-supported justification that is based on easily accessible sources.

Achievement: An agentic RAG pipeline that retrieves pertinent evidence by querying Cellosaurus, KEGG, PubMed, and the underlying omics data. Then, a language model is forced to synthesize arguments that only reference these recovered sources—never speculating or creating new evidence. Justifications are guaranteed to be based on actual data and literature thanks to this design.

{\bf Aim 5: Deliver an API and a User Interface with Exportable Reports.}

Goal: Make the system usable both visually (for human exploration) and programmatically (for incorporation into scientific operations).

Accomplishment: A stateless JSON API with four routers (genes, ranking, cell lines, and filters) is exposed using a FastAPI backend. SearchBar, FilterPanel, ResultsTable, and DetailPanel are components of a React web frontend. Results can be exported by users as PDF, CSV, or JSON. This system is run on a single Compute Engine instance (e2-standard-2) in zone asia-south2-b of Google Cloud Platform. The deployment is done completely by hand via SSH using the commands git pull, docker-compose up -d --build for the backend and Neo4j containers, and npm run build and nginx reload for any changes to the frontend. Certbot provides HTTPS.

In addition to these five specific objectives, the project included CRISPR gene dependence data and a real multi-hop knowledge graph using Neo4j, both of which were included in the initial brief and verified outcomes using a separate set of test genes provided by the industry partner AstraZeneca.

\section{Research Questions and Dissertation Structure}

Three research questions drove the project. Can multiple heterogeneous omics sources be integrated into a single reliable ranking despite incompatible identifiers and units, without discarding partially-covered cell lines? Does differentiating genes by biological class improve ranking quality over a single global scheme and if so, by how much and for which classes? Can automated, citation-backed justifications make such rankings interpretable and trustworthy to domain scientists without the justifications introducing unsupported or fabricated claims?

Chapter 2 provides technical and biological background: cell lines and gene expression, the five integrated data sources, why harmonisation is non-trivial and a critical review of related tools. Chapter 3 describes system design and implementation: data pipeline, the classical statistical model, gene-class aware weighting, knowledge graph integration, agentic AI model, supporting components and system architecture including UI and deployment. Chapter 4 critically evaluates accuracy, reliability and performance: functional testing, three critical bugs identified and corrected, ranking accuracy and the gene-class hypothesis, domain-expert validation including the discovery of a real TP53 defect, a systematic investigation into pathway-coherence scoring (4 methods tested, all failed, mechanistic explanation provided), method comparison and explicit limitations. Chapter 5 concludes with achievement against aims, key scientific findings and prioritised future work.

% =============================================================================

\chapter{Background}
\label{chap:background}

\section{Cell Lines in Cancer Biology Research}

Cell lines are populations of cells adapted to grow indefinitely in laboratory culture, derived from tumour biopsies, bone marrow or other tissue sources. They are the workhorse of cancer biology and drug development, enabling systematic study of cancer-specific properties in controlled environments. Cellosaurus \cite{bairoch2018cellosaurus} catalogues these lines, assigning stable accessions and recording tissue of origin, disease association, known synonyms across databases and clinical history. The same physical cell line is referred to by different names across different databases: HCC827, Hcc-827 and DepMap's ACH-000164 all denote the same lung adenocarcinoma line with EGFR exon 19 deletion. One of the main requirements was to resolve each reference onto a specific Cellosaurus accession, such as CVCL\_0063. In the absence of this, evidence would be fragmented and meaningful integration would be impeded since the same cell line would be treated as several different entities.

In cancer research, cell lines hold a special place. They lack the cellular heterogeneity and microenvironment of primary tumours, making results from cell lines sometimes non-predictive of patient responses. Yet they offer reproducibility, scalability and experimental tractability that in vivo models cannot match. A well-chosen cell line recapitulates key properties of the disease being studied: EGFR-addicted lung cancer models (HCC827, H3255) harbour activating EGFR mutations and depend on EGFR signalling for survival. TP53-deficient models (HCT116 with TP53 knockout, H1299) provide systems to study tumour suppressor loss. Triple-negative breast cancer models (MDA-MB-231, HCC1937) lack oestrogen, progesterone and HER2 receptors, matching a distinct clinical subtype. Selecting the wrong cell line can produce misleading results. A study of EGFR inhibitors run on EGFR-wild-type cells will show no effect, producing false-negative results.

The challenge of cell line selection has grown more acute as the number of available lines has expanded. Decades ago, researchers might choose from 50-100 available cell lines. Today, over 2,076 cancer cell lines are catalogued in public databases, each with different genetic backgrounds, tissue origins and molecular properties. Manual selection from such a large space is increasingly impractical. A scientist selecting a cell line for a study involving EGFR and KRAS mutations faces a combinatorial problem: check which cell lines express EGFR at appropriate levels, which have KRAS mutations, which are appropriate for the target tissue, which are proven to work in similar experiments.

\section{Gene Expression and Molecular Characterisation}

Gene expression varies significantly by cell type, disease state and even between individual cells of the same type. RNA-sequencing produces normalised read counts expressed as transcripts per million (TPM), a measure of relative abundance. Mass-spectrometry proteomics measures protein abundance directly via ionisation and detection of peptides, reporting absolute or relative quantities on incomparable scales. RNA and protein evidence are complementary rather than interchangeable. A gene may be transcribed at high levels but post-translationally degraded, resulting in low protein. Conversely, a gene may be transcribed modestly but translated efficiently and stabilised, resulting in high protein. Ranking cell lines by RNA expression alone misses protein-level effects.

Raw expression level should not be used to rank every gene. Because of their strong, tissue-specific, and cancer-type-specific expression, receptor tyrosine kinases like EGFR and ERBB2 are researched. Lung adenocarcinomas overexpress EGFR, while certain breast malignancies overexpress ERBB2. In these situations, pharmacological sensitivity and functional reliance are correlated with expression. Models where EGFR is implicated and inhibitors are probably effective can be accurately identified by ranking cell lines according to EGFR expression. On the other hand, because housekeeping genes like GAPDH or ACTB encode proteins necessary for fundamental cellular function, they are expressed at comparable quantities in almost all cell types. When choosing a cell line, their expression is constitutive and unhelpful. Since GAPDH is expressed everyplace, searching for GAPDH shouldn't provide "high-GAPDH" cell lines. Instead, it ought to provide cell lines suitable for the experimental setting.

Genes that have lost their ability to function, such as TP53, BRCA1, or PTEN, fall into a completely separate category. These genes are studied not because they are overexpressed but because they are absent or mutated. TP53 is a tumour suppressor: wild-type TP53 is expressed in nearly all cell lines at baseline levels. Cell lines with wild-type, functional TP53 are not appropriate models for studying tumour suppressor loss. Cell lines with TP53 mutations or deletions are the scientifically relevant models. For these genes, expression level is nearly irrelevant; the relevant signal is mutation status and impact. A model for EGFR-addicted lung cancer and a model for TP53-deficient cancer require completely different selection criteria. Expression ranking works for the first and fails for the second. A single global weighting formula applied uniformly to all genes cannot capture these distinctions.

\section{The Five Public Data Sources}

CellSelector Omics integrates five public databases, each covering different aspects of cell biology with different strengths and limitations.

{\bf Human Protein Atlas (HPA).} The largest unified transcriptomics dataset for cancer cell lines is the HPA \cite{uhlen2015proteome}, which offers RNA-seq normalized expression across 1,198 cancer cell lines. Transcripts per million, or TPM, is a standardized metric that allows for direct comparison of HPA expression data. For the 1,198 cell lines that are present, coverage is quite high; yet, compared to all known cancer cell lines, coverage is insufficient. Because each cell line was treated in the same way, batch effects that affect diverse datasets were eliminated, which is one of HPA's strengths. 

{\bf Cancer Dependency Map (DepMap).} Log-transformed TPM expression and CRISPR gene dependence scores for 1,399 cell lines are provided by DepMap \cite{tsherniak2017depmap}. DepMap expression is reported as $\log_2(\text{TPM}+1)$, a log-transformed scale that amplifies differences at lower expression levels and dampens differences at high expression. This is different from HPA's linear TPM, requiring unit conversion. DepMap's unique contribution is CRISPR dependency scores: the probability that knocking out each gene is lethal to each cell line. For loss-of-function studies, CRISPR dependency is highly informative. A cell line with high CRISPR dependency on TP53 is TP53-functional and sensitive to TP53 loss; a cell line with low dependency is TP53-mutant or otherwise insensitive. DepMap covers more cell lines than HPA (1,399 vs. 1,198) but lower expression standardisation. DepMap is the primary data source for dependency studies.

{\bf Gene Expression Omnibus (GEO).} GEO \cite{barrett2013geo} aggregates heterogeneous third-party submissions from 3,159 datasets, representing hundreds of thousands of independent experimental studies. GEO is vast but highly heterogeneous. Different experiments use different measurement protocols (microarray, RNA-seq), different normalisation methods, different tissues and disease states. A single GEO dataset covers 10-100 cell lines; no single experiment covers all 2,076. GEO expression values are unitless, with original normalisation unknown. The strength of GEO is comprehensiveness and independence; these data come from researchers worldwide, not from a single curated source. The flaw is heterogeneity; batch effects and technique variations might lead to inaccurate results when directly averaging among GEO datasets. For CellSelector Omics, GEO serves as a confirmatory source, adding evidence when available but not required for any recommendation.

{\bf CCLE Proteomics.} CCLE \cite{nusinow2020proteomics} provides mass-spectrometry protein abundance across 375 cell lines, the only direct protein-level measurement in this project. Protein abundance and transcript abundance are only moderately correlated ($r \approx 0.4$--$0.6$); a gene can be transcribed without translation and translated without transcription. CCLE proteomics is therefore essential for capturing post-translational aspects of cell biology. The limitation is sparse coverage (375 of 2,076 cell lines, 18.1\%). For cell lines without CCLE data, protein-level evidence is unavailable. This sparsity is a known limitation acknowledged throughout the project.

{\bf KEGG Pathways.} KEGG \cite{kegg2016} provides manual curation of biological pathways linking genes to cellular processes. The ErbB signalling pathway includes 87 genes involved in receptor tyrosine kinase signalling. The TP53 pathway includes genes involved in DNA damage response and apoptosis. Pathway membership enables graph-based ranking: cell lines where genes in the same pathway show coordinated activity might be appropriate for studying that pathway. However, as Chapter 4 discusses, transcriptomics-based pathway detection has fundamental limitations, particularly for post-translational activation mechanisms.

\section{The Data Integration Challenge}

Why is integrating these five sources non-trivial? The answer lies in incompatible identifiers, units and coverage patterns.

{\bf Identifier Mismatch.} Cell lines are identified by Cellosaurus accession (CVCL\_0063), DepMap's ach\_id (ACH-000164), or free-text names (HCC827, H3255, various spellings and synonyms). Genes can be identified using RefSeq accession, gene symbol such as EGFR or Ensembl gene ID, like ENSG00000005471. The simple use of string concatenation leads to inconsistencies. For example "EGFR" could be used in one database while " growth factor receptor" in another – both names are for the same gene although string matching fails here. Or even if "HCC827" and "Hcc-827" represent the cell line they will differ from each other due to case-sensitive search. The HLA class of genes is a case when "HLA-A" and "HLA*A02:01" correspond to the same gene but, with different specificity. The CellSelector Omics uses the Cellosaurus and Ensembl to resolve any cell line names whereas all the genes are treated as canonical IDs.

{\bf Unit Incompatibility.} DepMap gives values in $\log_2(\text{TPM}+1)$ HPA gives linear TPM and GEO gives values. Even though the numbers look similar they mean things. A value of 10 in the $\log_2$ scale means $2^{10} = 1{,}024$ TPM in terms. But 10 TPM in DepMap is not the same as 10 TPM in a scale—it's about a thousand times different. If you average these directly you get a number that doesn't make sense. Converting from $\log_2$ to linear is easy mathematically. It’s not enough. Values on a shared linear scale still depend heavily on the sample used. A group of cells with expression will push the distribution up. A group with expression will pull it down. So a cell line might seem expressed in one dataset just because that experiment happened to use cells, with strong gene expression.

{\bf Normalisation Strategies.} Two complementary normalisation approaches address these challenges. Z‑score normalisation converts each value to (value minus mean) divided by \texttt{standard\_deviation} standardising against the source’s distribution. Z‑score normalisation preserves the magnitude of outliers. Assumes comparable distribution shapes across sources. Percentile normalisation converts each value to its rank within the source discarding absolute magnitude but remaining invariant to monotonic transformation. A cell line ranking at the percentile in HPA is treated exactly the same as one at the 90th percentile in DepMap no matter what the absolute TPM values are. The project uses both Z‑score normalisation and Percentile normalisation fused through rank fusion (RRF). RRF dampens the influence of ranks and produces a single score, from multiple rank lists.

{\bf Coverage Patterns.} Different subsets of cell lines are covered by various sources. HPA covers 1,198, DepMap covers 1,399, GEO covers thousands across datasets, CCLE covers 375, KEGG covers all 2,076 (pathway membership is inferred). A naive integration strategy requiring complete coverage across all sources would discard any cell line missing data from even one source. If a cell line is in HPA but not DepMap, it is dropped. If present in DepMap but not HPA, it is dropped. This discards the majority of cell lines. Alternatively, CellSelector Omics averages scores across available sources and tracks coverage separately. A cell line present in HPA and DepMap but not CCLE is ranked based on the two available sources, not penalised for missing the third. This preserves 2,076 cell lines rather than reducing to a small subset with complete coverage. The trade-off is that recommendations for sparsely-covered cell lines rest on less evidence and are weighted as such.

\section{Related Work and Existing Tools}

Multiple tools exist for cell-line characterisation, each with distinct purposes and limitations.

{\bf DepMap Interactive Browser.} Rich interactive dependency and expression data visualization is provided by DepMap \cite{tsherniak2017depmap}. Researchers may look up a cell line and see its expression profile, dependence scores, and thousands of genes' worth of genomic features. The tool is powerful for hypothesis exploration: "Given that I know cell line X, what genes is it dependent on?" DepMap excels at answering this question. However, DepMap does not rank cell lines for prospective experiments. It does not answer "I want a cell line expressing gene Y with mutation Z; what are the top candidates?" To answer such questions, scientists must manually cross-reference DepMap with other sources, an integration burden that falls entirely on the user.

{\bf cBioPortal.} cBioPortal \cite{gao2013cbioportal} visualises cancer genomics within specific studies. Scientists can explore which mutations co-occur in tumours, which genes are altered in patient samples, how alterations correlate with outcomes. cBioPortal is designed for hypothesis exploration within a study, not for cell-line model selection. It visualises patient data, not cell-line panels. The integration burden again falls on the user.

{\bf CELLector.} CELLector \cite{iorio2016cellector} matches genomic features (mutations, copy-number alterations, gene expression) to patient tumour subpopulations, returning cell-line sets that phenocopy specific patient cohorts. CELLector answers "Which cell lines best represent my patient population?" rather than "Which cell lines express my gene of interest?" The tool returns matched sets rather than ranked lists and does not estimate confidence in matches.

{\bf Cell Model Passports.} Cell Model Passports catalogue clinical and functional annotation for cell lines, including tissue of origin, disease, growth medium and doubling time. The resource is comprehensive but browsed rather than queried programmatically. Searching for cell lines requires manual filtering.

{\bf Common Limitations.} All existing tools have a common limitation: the integration burden falls entirely on the user. A scientist must visit DepMap, HPA, GEO, COSMIC and CELLector separately, manually cross-reference results, mentally integrate evidence and make a decision. This is time-consuming, error-prone and produces inconsistent results. Different researchers checking the same criteria may choose different lines based on which sources they consulted, in which order and what they happened to find. There is no canonical, evidence-backed, reproducible recommendation. CellSelector Omics automates this integration, moving from browsing multiple sources to ranked, justified recommendation in a single interface. The system combines evidence from all sources, weights it according to gene class, surfaces confidence and produces publication-ready justifications.

\section{The Research Gap}

Although it is still ad hoc and manual, cell-line selection is a fundamental stage in experimental design. There is currently no tool that offers ranked recommendations with stated confidence for a multi-gene query (e.g., "I want a cell line expressing EGFR, with KRAS mutation, lacking TP53"). This integration issue is one that scientists must resolve on their own. By creating a system that combines diverse data sources, employs principled ranking logic, and presents arguments supported by evidence, this project fills this gap. Speed and repeatability are the outcomes. The recommendations are the same whether the same query is run on various days. A scientist can provide detailed data and references to relevant literature to support their decision. It now just takes minutes to complete the selecting procedure, which used to take hours.

% =============================================================================

\chapter{System Design and Implementation}
\label{chap:execution}

\section{Research Methodology and Design Philosophy}

CellSelector Omics was created based on three ideas that aimed to solve problems in other tools.

{\bf Principle 1: Data-Centric Integration Without Loss.} We tried to avoid discarding cell lines that lacked any data from any source and instead integrated in a manner that would not punish missing coverage. If more than one cell line had a percentile rank score, we averaged them, not added them together, and did not deduct any from the cell line that did not have a third score. That is a direct response to a common weakness in current tools: that of incomplete data coverage, which would exclude otherwise good candidates.

{\bf Principle 2: Interpretability-First Design.} Every suggestion has to be clear and able to be followed to the original sources. This meant we said no to black-box machine learning even if it could be more accurate. We used a clear scoring system where each type of evidence (like RNA levels, protein amounts, mutation details and pathway information) adds a named part with a weight to the final score. Agentic AI uses retrieval-augmented generation to make explanations that are backed by references so scientists can see where the suggestions come from. A tool that helps make decisions is only useful if the person using it understands and believes in the reasons behind the choices.


{\bf Principle 3: Biologically-Aware Ranking.} Different types of genes need ways to be ranked. Genes that are only active in tissues do well when looking at expression levels but genes that are active everywhere don't. Genes that lose their function do best when looking at mutations. A single formula, for all genes can't handle these differences. This led to the idea that genes should be grouped and ranked in ways. This idea was. Proved in Chapter 4. The idea came from seeing that an early model gave results that didn't make sense for some types of genes.

\begin{figure}[t]
\centering
\includegraphics[width=0.9\textwidth]{figures/pipeline.png}
\caption{Emphasises harmonisation as cohesive workflow, contextualises error rates and coverage bottlenecks, flows naturally toward knowledge graph construction}
\label{fig:data-schema}
    \end{figure}


\section{Data Pipeline and Harmonisation}

The pipeline reads 222 million expression records and 21.5 million CRISPR dependency records from five sources. The total size of the input files is 1.1 gigabytes. The pipeline streams data in chunks of about 2,000 gene columns adding them to Parquet storage. The maximum memory used by the pipeline never exceeded 500 megabytes so the entire import could run on a laptop.

Every input record is linked to a Cellosaurus identifier by normalising cell line names and synonyms against the Cellosaurus database. Expression records are ranked by percentile inside each source to adjust for varying units and scales.

Four preprocessing challenges were solved. First linking records across sources needed the pipeline to use DepMaps \texttt{ach\_id} as the key. Second the pipeline had to convert DepMaps $\log_2(\text{TPM}+1)$ values to ordinary TPM so that units matched. Third the pipeline had to find and clean invalid values, such as N/A and infinite numbers and eliminate duplicate rows. Fourth the pipeline changed the data shape from wide to long format using the pair (\texttt{ach\_id}, \texttt{ensembl\_id}) as the key, which made later filtering easier.

\section{The Classical Statistical Model}

The algorithm scores cell lines in four steps when ranking a gene.

{\bf Step 1: Per-Source Ranking.}In each source, the raw expression value of a cell line is transformed to a percentile rank, which is the percentage of candidate cell lines with expression values as low or lower. It ensures that all sources are compared on a 0 to 1 scale which is independent of the natural range of the gene expression, allowing for a fair comparison of the relative position of a cell line among different sources for different genes and platforms.

{\bf Step 2: Fusing Ranks via Reciprocal Rank Fusion.} Each ranking is turned into a weight using RRF \cite{cormack2009rrf}. RRF reduces the impact of ranks. For each cell line up to six ranking lists are combined by averaging, not by adding them up. Averaging makes sure that a cell line that does not have data from every source is not unfairly scored lower than a cell line that has all the data. The amount of coverage for each source is recorded separately. Coverage is tracked separately and filtered out later rather than being built into the score..

{\bf Step 3: Protein Ranking.} Protein intensity is handled the same way. Two ranking lists are fused because protein data comes from one source.

{\bf Step 4: Combining into Final Score.} After the RNA and protein RRF scores are fused they are min‑max normalised to a range of 0 to 1. Then they are combined, using default weights of 0.7 for RNA and 0.3 for protein. A cell line can be in one of four situations: both RNA and protein contribute, only RNA, only protein or neither.

When we ask about genes at once each gene is ranked on its own. After that the ranks for each gene are fused into one ranking.

\section{Gene-Class Aware Weighting}

Not every gene should be ranked by expression. That idea led us to use weights for each class of gene.

Tissue‑specific genes are expressed a lot and only in cancer types. For these genes ranking by expression works well.

Ubiquitous genes are expressed at the same level, in almost every cell line. Ranking by expression gives no difference.

Loss‑of‑function genes are studied by looking at their absence or mutation not their expression levels. Ranking by expression tells us nothing; we need to look at mutation and context.

One single weighting formula cannot handle all these differences.

We tuned the weights using leave‑one‑out cross‑validation on a set of 44 genes aiming to maximise Mean Reciprocal Rank. For each gene that we left out we refitted the weights using the other genes in the same class before scoring which helps avoid overfitting.

\begin{table}[t]
\centering
\begin{tabular}{|l|c|c|c|c|}
\hline
{\bf Gene Class} & {\bf RNA} & {\bf Protein} & {\bf Quality} & {\bf Mutation-Impact} \\
\hline
Tissue-specific & 0.60 & 0.10 & 0.20 & 0.10 \\
Ubiquitous & 0.50 & 0.05 & 0.35 & 0.10 \\
Loss-of-function & 0.003 & 0.048 & 0.046 & 0.80 \\
\hline
\end{tabular}
\caption{Optimized evidence weights per gene class, prior to the later integration of graph-based (RWR) scoring discussed in Section 4. The optimizer, constrained only by ranking performance with no biological instruction, converged on markedly different weightings. For loss-of-function genes, expression is nearly discounted (0.003) in favour of mutation-impact scoring (0.80). These weights emerge from data rather than design, providing the strongest evidence for the gene-class hypothesis.}
\label{tab:learned-weights}
\end{table}

The most striking finding: for loss-of-function genes, the optimizer independently assigned 0.80 weight to mutation-impact evidence and only 0.003 to expression. A biologist would state this principle immediately. Yet it emerges purely from ranking performance. This validates the gene-class distinction empirically rather than by assumption.

\section{Random Walk with Knowledge Graph with Restart}

A knowledge graph with hops was built in Neo4j connecting 19,000 genes, 255 KEGG pathways, and 2,076 cell lines. The edges in this graph represent expression relationships, mutation interactions and pathway membership. To measure how close a cell line is to a query gene, Random Walk with Restart (RWR) was used This method relies on PageRank to compute structural proximity within the graph. The RWR scores are added as a layer of evidence alongside traditional scores. Weights specific, to RWR were chosen using grid search to performance.

\begin{figure}[t]
\centering
\includegraphics[width=0.9\textwidth]{figures/rwr_mechanism.png}
\caption{Stratified by gene class with functional context; explains why each class benefits differently from network integration; stronger technical language ("under-ranking by expression-only methods")}
\label{fig:rwr-mechanism}
\end{figure}


\section{Agentic AI Model with Citation-Backed Justifications}

The agentic component uses retrieval-augmented generation (RAG) to produce structured justifications for recommendations. When a user requests justification for a specific result, the system queries PubMed, KEGG, Cellosaurus and underlying omics data to retrieve relevant evidence, then uses a language model with explicit instructions to synthesise citation-backed explanations. Justifications are grounded in retrievable sources rather than speculative.

Design principle: ranking decides, language model explains. The ranking produces scores independently. The language model cites only from retrieved sources, never inventing evidence or speculating. Any hallucination in justification does not alter the underlying recommendation.

\section{Supporting Components}

The similarity component of the ranking attempts to answer a different question than the main ranking: not "What cell lines do I expect to see that gene" but "What cell line do I expect this one to look like". Similarity is calculated over a 19-dimensional feature vector combining expression evidence (RNA, protein, quality, and context scores where observed) with cell-line-intrinsic genomic features - microsatellite instability (MSI) score, chromosomal instability index, and an evidence-completion measure indicating for each cell line the fraction of sources providing supporting evidence. Each feature is normalised independently to the 0-1 interval - the median of missing values is used for each feature individually (rather than a colum-wise median across features). Similarity is the cosine of the angle between the weighted feature vectors, which is exactly the same normalisation and distance metric as is used in the model's own ranking step - ensuring a conceptually consistent notion of "cell line more like" and "more well ranked".

\begin{table}[t]
\centering
\begin{tabular}{|l|c|c|c|}
\hline
{\bf Feature set} & {\bf Same} & {\bf Different} & {\bf Gap} \\
\hline
Expression, min-max normalised & 0.650 & 0.548 & +0.102 \\
\hline
\end{tabular}
\caption{Mean cosine similarity within and between lineages, computed using the min-max normalised, 19-dimensional feature matrix (EGFR query, $n = 15,238$ total sampled pairs). Same-lineage cell lines show a meaningfully higher mean similarity than different-lineage cell lines, confirming the feature representation captures lineage-relevant structure.}
\label{tab:similarity-scaling}
\end{table}

\section{System Architecture and Deployment}

{\bf Data Pipeline.} Fourteen datasets pass through a three-stage offline pipeline. Python scripts clean each: cell line names resolve to DepMap accessions, gene symbols to Ensembl identifiers, DepMap's $\log_2(\text{TPM}+1)$ values convert to linear TPM, wide matrices pivot to long format. Each CSV converts to Apache Parquet using PyArrow with Snappy compression and dictionary encoding in 500,000-row chunks. At startup the backend registers each Parquet file using pandas, so the data-access layer queries via standard SQL.



{\bf Backend API.} The backend is a Python web service built with FastAPI \cite{fastapi2021}, exposing a stateless JSON API. The gene router provides prefix-based autocomplete and symbol-to-Ensembl resolution. The ranking router accepts POST bodies specifying target genes with expression directions, weights, filter modes and exclusion criteria, delegating to the ranking service. The cell line router serves evidence-breakdown, cosine-similarity and comparison endpoints. The filter router returns disease, lineage and subtype values with valid pairings for cascading dropdowns.

{\bf Frontend User Interface.} The frontend is a single page app built with React \cite{react2024}. Packed with Vite \cite{vite2024}. I like how the app talks to the backend using a JSON API. All rendering state lives inside React hooks. SearchBar offers debounced autocomplete. Every added gene gets a high/low toggle. FilterPanel shows cascading dropdowns and a filter/soft-boost toggle. The advanced section of FilterPanel lets you choose mutation/fusion modes, toggle RNA sources slide the RNA:protein weight, set instability thresholds and exclude metabolites or miRNA. ResultsTable shows ranked results in a table with dynamic columns and a CSV export button. DetailPanel slides in as an overlay and displays per-source scores, protein data, mutation annotations and the five similar cell lines. CompareView opens a modal that lets you compare rows side by side with tabs for each gene when you query genes.

React component architecture builds a modular hierarchy that keeps responsibilities separate. SearchBar handles debounced autocomplete with real-time validation. FilterPanel keeps track of cascading state. ResultsTable can render over 1,500 rows by using virtualisation. DetailPanel and CompareView are overlays that do not block scrolling of the base results table. State management relies on React hooks avoiding external libraries, which keeps dependencies and bundle size small. Production builds aim, for a page load time under 200 milliseconds.

{\bf Deployment.} The application is deployed to a single Virtual Machine (VM) instance of Google Cloud Platform Compute Engine with 2 vCPUs and asia-south2-b zone (e2-standard-2). The back-end is implemented as a Docker container using docker-compose and serves the FastAPI application using Uvicorn. This knowledge graph is hosted in a separate self-hosted Neo4j Docker container (cso-neo4j) on the same VM. It's an offline build of the frontend (npm run build) and forwarded via Nginx running on the same VM instance, not in a separate container or through CDN or Cloud Storage. The harmonised expression and proteomics information is stored in parquet files in the VM local disk. No automated CI/CD was put in place, all the deployment is carried out manually via SSH on the VM running the commands git pull, docker-compose up -d --build to rebuild and launch the backend container, npm run build and nginx reload to update the frontend. The deployment time varies depending on the changes needed and will take a couple of minutes because system is rebuilt from scratch during deployment, but downtime has been experienced during production incidents (e.g., 971s of downtime due to the loss of in-memory cache) to prove that manual deployment without automated testing can have operational risk. The focus was simplicity, and not automation: a manual approach worked for a production research tool, but it was fragile.

\begin{figure}[t]
    \centering
    \includegraphics[width=0.9\textwidth]{figures/cellselector_system_architecture.png}
    \caption{System architecture showing data pipeline (left), GCP Compute Engine deployment with FastAPI backend, Neo4j knowledge graph, and Nginx-served React frontend (center and right), and agentic justification pipeline (bottom). All components run on a single e2-standard-2 VM with manual SSH-based deployment.}
    \label{fig:system_architecture}
\end{figure}



% =============================================================================

\chapter{Critical Evaluation}
\label{chap:evaluation}

There is no ground‑truth ranking in this problem space. Researchers have made lists of cell‑line associations. No outside gold standard tells which recommendations are best for any gene. Evaluation proceeds along four axes: functional correctness, ranking validity, methodological rigour and domain‑expert validation. Each axis asks a question, about evaluation: does evaluation compute results? Does ranking align with known biology? Are methods ? Does a real scientist find evaluation results sensible? We validate against literature-curated associations, test genes provided by AstraZeneca, solicit independent domain-expert predictions and systematically investigate failure modes to extract mechanistic insight.

\section{Evaluation Strategy}

A 44-gene validation set stratified by biological class provides the primary evaluation benchmark. This set was constructed by identifying genes with well-established, literature-documented associations to specific cell lines, then partitioning genes by the biological class hypothesis described in Chapter 3.

{\bf Tissue-specific genes (n=24):} Genes that are high and specifically expressed in particular cancer types and are studied because of this tissue/cancer-type-specific expression. Examples include EGFR in lung adenocarcinoma, AR in prostate cancer, KIT in leukaemias, PGR in oestrogen-receptor-positive breast cancer. For these genes, expression-based ranking should correctly identify appropriate models. These genes are studied because overexpression or gain-of-function drives cancer and cell lines reflecting this biology are appropriate experimental models.

{\bf Ubiquitous genes (n=4):} Genes expressed at broadly similar levels across nearly all cell lines because they encode proteins essential for basic cellular function. Examples include TP53, GAPDH, PARP1, CCND1. These genes are typically not studied by selecting high-expressing cell lines; GAPDH is used as a normalisation control precisely because its expression is uninformative for distinguishing cell lines. TP53 is studied through absence or mutation, not expression. Expression-based ranking offers no discriminative signal for these genes and can produce biologically nonsensical results.

{\bf Loss-of-function genes (n=16):} Genes studied through absence or mutation rather than expression, including canonical tumour suppressors BRCA1, BRCA2, RB1, TP53, PTEN and others. For these genes, the relevant signal is mutation status and functional impact, not expression level. Cell lines with damaging mutations are appropriate models; cell lines with wild-type, functional copies are not.

Of 62 known canonical cell lines documented in the literature for these genes, 60 were successfully resolved to Cellosaurus identifiers and were present in the unified data store (96.8\% coverage). The two unresolvable lines were historical names with multiple competing modern identifiers or unclear provenance.

\section{Functional Testing}

{\bf Base-Case Query: EGFR High Expression.} I ran a query for EGFR with an inclusion directive and no filters. The system returned all 2,076 ranked cell lines with fields: Cellosaurus accession, tissue of origin, disease, overall ranking score, per-source evidence breakdown, confidence percentile and coverage flags indicating which data sources contributed. All scores lay within the range [0 1]. Were sorted in descending order. No null or NaN values appeared in the output. The top‑ranked result was HCC827, well documented as an EGFR‑amplified exon‑19‑deletion‑mutant lung adenocarcinoma cell line with known hypersensitivity, to EGFR tyrosine kinase inhibitors. This result is biologically correct validating that base‑case queries produce output.

{\bf Multi-Gene Query: EGFR + SEC61G.} Tested a query combining two genes: EGFR (HIGH) and SEC61G (HIGH). The system returned ranked results with a breakdown showing how each gene contributed to the ranking. Both genes are located close to each other on chromosome 7p11.2 and SEC61G is often amplified alongside EGFR, in lung cancers. The system accurately identified cell lines with expression of both genes.

{\bf Cumulative Filtering: Disease, Tissue, Lineage.} I tested filtering on EGFR queries. When I did not filter the system returned all 2,076 cell lines. When I filtered by lung cancer I got 412 results and every one of those 412 was confirmed to be lung cancer. When I filtered by tissue and subtype small cell lung cancer I found 18 results and each of those 18 was confirmed to be small cell lung cancer. The cell line HCC827 appeared in both the results and the lung cancer filtered results and HCC827 had the same score in both. HCC827 was correctly missing, from the small cell lung cancer filtered results because HCC827 is an adenocarcinoma and not small cell lung cancer.

{\bf Mutation Filtering: Include/Exclude/Ignore Modes.} Testing three modes on EGFR. Ignore (default): 2,076 results returned. Include (any EGFR mutation): 1,247 results, which's 60.1 percent of all cell lines. Exclude (no EGFR mutation): 829 results, which's 39.9 percent of all cell lines. Verification: 1,247 plus 829 equals 2,076 with no entries. Importantly the ranking score for a cell line stayed the same whether it appeared in include mode or ignore mode. This proves that the EGFR mutation filter works as a filter not as a penalty, to the score.

\section{Rigour of Development: Critical Bugs Identified and Corrected}

Three non-trivial correctness bugs were identified and corrected through systematic testing and cross-validation rather than incidental discovery. Each bug had potential to silently invalidate downstream evaluation results and each was discovered only through deliberate, meticulous checking rather than obvious failure. We report these in full as direct evidence of methodological rigour, not merely as changelog.

{\bf Bug 1: HPA and GEO Identifier Misattribution.} While investigating an apparent hallucination in the agentic justification pipeline--generated rationales confidently cited specific HPA expression values for cell lines later found to have no HPA coverage at all--we traced the fault to the load\_mappings() function. This function constructed HPA and GEO lookup dictionaries by calling .dropna() independently on two columns of a lookup table before zip()-ing them together. Because the two columns (cell-line names and Cellosaurus identifiers) contained different numbers of missing values, this silently produced positionally misaligned key-value pairs. Each HPA name or GEO accession was paired with the n-th non-null Cellosaurus identifier by row position, not with its own row's correct identifier.

We quantified the resulting damage precisely. Of 1,198 HPA mappings, 1,126 (94.0\%) pointed to the wrong cell line. Of 3,159 GEO mappings, 2,534 (80.2\%) pointed to the wrong cell line. The DepMap mapping, constructed via itertuples() rather than a paired dropna operation, was unaffected. This corruption sat upstream of every RNA-based evidence score in the entire pipeline. Correction required filtering both columns jointly with dropna(subset=[...]) on a single dataframe before pairing them, preserving row alignment throughout.

Because this function fed every RNA-based evidence score, its correction required re-running the full evaluation chain from scratch. We explicitly discarded several earlier findings--including an apparently clean, monotonic pathway-weighting result for loss-of-function genes--that we subsequently determined had been artefacts of this corrupted mapping. All evaluation results reported in this dissertation reflect the corrected pipeline only.

{\bf Bug 2: Non-Deterministic Tie-Breaking.} Repeated evaluation runs under nominally identical configurations produced measurably different Mean Reciprocal Rank values. Across three successive runs of the same configuration, we observed MRR = 0.2421, then 0.2121, then 0.1855--a far larger variation than floating-point rounding error would explain. We traced this to two independent sources of non-determinism. First, Python's hash-seed-randomised set iteration order altered the summation order of floating-point reciprocal ranks by small but non-zero amounts across runs. Second, unstable sort operations wherever candidate cell lines shared identical scores caused the specific cell line selected as ``top-ranked'' among a tied group to vary between otherwise identical runs.

We resolved both sources by enforcing sorted() iteration over any set consumed by an aggregation step and by applying stable, multi-key sorts (kind=``stable'') at every point where scores feed into rank computation. We verified the fix by running the same configuration three times in separate processes with randomised PYTHONHASHSEED, confirming bit-identical MRR of 0.12403927790842631 across all three runs, demonstrating that non-determinism was completely eliminated.

{\bf Bug 3: Score-Saturation Tie-Break Artefact.} Following the introduction of mutation-impact scoring for loss-of-function genes, an initial cross-validation run produced a striking result: HCC1937, the canonical BRCA1-deficient reference line, achieved held-out reciprocal rank of exactly 1.000. On detailed inspection, we discovered an unintended interaction between the scoring formula and its clipping step. For cell lines carrying a maximally-scored mutation (mutation\_impact\_score = 1.0), the weighted sum of mutation, GEO-confirmation and residual expression terms exceeded 1.0 before clipping. The clipping step then erased precisely the fine-grained differences intended to break ties among such highly-scored lines. As a result, roughly three cell lines were reduced to an identical maximal score of 1.000 before correction, with the deterministic tie-break subsequently ordering HCC1937 first based on its higher raw mutation-impact and RNA-expression scores relative to the other tied lines.

We corrected this not by altering the scoring formula itself--which had been validated independently via weight optimisation--but by introducing an explicit, deterministic multi-key sort applied wherever final scores feed into ranking. The sort proceeds as: primary sort by final score, ties broken by unclipped mutation score, then by RNA expression score, then by Cellosaurus identifier as a final deterministic fallback. Under this correction, HCC1937's true, tie-break-robust held-out reciprocal rank was 0.500 (rank 2 of the tied block), a materially more conservative and honest figure than the pre-correction value of 1.000.

We regard this third bug as particularly instructive methodologically. It demonstrates that a scoring formula can be internally correct and fully validated in isolation, yet still produce artefactual results once combined with naive ranking. This is a category of error that unit-level correctness testing would not have surfaced; we identified it only because we deliberately audited the tie structure of our best-performing individual result rather than accepting it uncritically.

\section{Ranking Accuracy and the Gene-Class Hypothesis}

Mean Reciprocal Rank (MRR) was the primary accuracy metric. MRR is the mean of the reciprocal of each gene's correct-answer rank: a rank-1 result contributes 1.000 to the mean, rank-2 contributes 0.500, rank-10 contributes 0.100. This metric penalises poor placement extremely steeply and is therefore dominated by a comparatively small number of genes for which no cell line in the dataset is a strong candidate. We caution against interpreting MRR as a percentage-style accuracy. For ubiquitous genes like GAPDH, no single best answer exists by construction, since the gene is expressed uniformly across hundreds of cell lines. For such genes, high MRR is unobtainable.

Weights were optimised via leave-one-out cross-validation (LOO-CV) on the 44-gene set to maximise MRR. For each held-out gene, weights were re-fitted on remaining genes of the same class before scoring, avoiding the optimistic bias of in-sample evaluation and the data-leakage problem of reusing test data during training.

\begin{table}[t]
\centering
\begin{tabular}{|l|c|c|c|c|}
\hline
{\bf Gene Class} & {\bf n} & {\bf Global Weighting} & {\bf Per-Class Weighting} & {\bf Improvement} \\
\hline
Tissue-specific & 24 & 0.1547 & 0.3300 & +113\% \\
Ubiquitous & 4 & 0.0062 & 0.1216 & +1,861\% \\
Loss-of-function & 16 & 0.0041 & 0.1968 & +4,700\% \\
\hline
{\bf Overall} & 44 & 0.1096 & 0.2367 & +116\% \\
\hline
\end{tabular}
\caption{Leave-one-out cross-validated Mean Reciprocal Rank by gene class. Per-class weighting improves overall MRR from 0.1096 to 0.2367 (116\% improvement). Loss-of-function genes show largest proportional improvement (4,700\%), consistent with the hypothesis that these genes are biologically determined by mutation status rather than expression. The improvements are stable and reproducible, not artefacts of overfitting.}
\label{tab:mrr-by-class}
\end{table}

{\bf Tissue-Specific Genes.} These genes achieve high accuracy under per-class weighting. Top performers include PGR (MRR 0.519), KIT (0.505), ALK (0.350) and EGFR (0.243). For all of these genes, the known canonical cell line ranked among the top results. The improvement from global (0.1547) to per-class (0.3300) weighting is 113\%, reflecting that expression-based ranking is fundamentally the right approach for this class.

{\bf Ubiquitous Genes.} These genes show improvement from 0.0062 to 0.1216, a 19.6-fold improvement. However, the absolute MRR remains modest. This is not a deficiency in the model; it reflects the underlying biology. For GAPDH, expressed uniformly across hundreds of cell lines, there is no biologically defensible ``best'' answer by construction. A scientist querying for GAPDH should not expect confident ranking; they should expect honest acknowledgement that GAPDH is not informative for cell-line selection and should choose their model on other grounds.

{\bf Loss-of-Function Genes.} These genes show the most dramatic transformation: from 0.0041 to 0.1968, a 48-fold improvement. Under global weighting, canonical TP53 loss-of-function models ranked in the hundreds or thousands--biologically nonsensical results reflecting that expression-based ranking measured the wrong signal entirely. Under per-class weighting informed by mutation status, they rank in the top 20. This improvement validates the gene-class hypothesis at the most extreme level.

{\bf Precision@k as Practical Utility Metric.} We caution against interpreting MRR as the only accuracy measure. MRR is dominated by genes for which no cell line is a strong candidate. With respect to practicality, what the scientist wants to know is whether the right cell line is included among a shortlist of potential candidates. This is where Precision@k (percentage of top-k matches which are right) comes into play. It is likely that the scientist will be able to examine a shortlist of 10 cell lines manually, meaning that Precision@10 would be the only practically relevant metric. With respect to our system incorporating ranking based on the graph model, Precision@10 significantly outperforms MRR.

\begin{table}[t]
\centering
\begin{tabular}{|l|c|c|c|c|c|}
\hline
{\bf Gene Class} & {\bf n} & {\bf P@1} & {\bf P@5} & {\bf P@10} & {\bf MRR} \\
\hline
Tissue-specific & 24 & 0.75 & 0.63 & 0.48 & 0.3300 \\
Ubiquitous & 4 & 0.25 & 0.15 & 0.10 & 0.1216 \\
Loss-of-function & 16 & 0.44 & 0.34 & 0.22 & 0.1968 \\
\hline
{\bf Overall} & {\bf 44} & {\bf 0.55} & {\bf 0.46} & {\bf 0.33} & {\bf 0.2367} \\
\hline
\end{tabular}
\caption{Precision@k across gene classes. Tissue-specific genes achieve high P@1 (0.75) 
and acceptable P@10 (0.48). Ubiquitous genes remain difficult; P@10 is low (0.10). 
Loss-of-function genes intermediate (P@10=0.22). Notably, Precision@10 across all 
queries (0.33) is substantially higher than MRR (0.2367), validating that results are 
actionable for practical shortlisting.}
\label{tab:precision-at-k}
\end{table}



\section{Comparison Against Alternative Ranking Methods}

To establish that our production scoring formula--a gene-class-adaptive weighted linear combination--represents a genuinely competitive choice rather than simply the first approach attempted, we benchmarked it against two established alternative ranking paradigms: Reciprocal Rank Fusion (RRF) and LambdaMART.

{\bf Reciprocal Rank Fusion.} RRF \cite{cormack2009rrf} combines evidence sources by their rank rather than raw score, dampening influence of low ranks via reciprocal-rank transformation. RRF is appealing because it is parameter-free (once evidence sources are chosen, no weighting decisions required) and robust to outliers and scaling differences.

{\bf LambdaMART.} LambdaMART \cite{burges2010} is a gradient-boosted learning-to-rank model capable of learning non-linear interactions between evidence features that a linear formula cannot express. If the relationship between evidence and correct ranking is fundamentally non-linear, LambdaMART should discover it.

This comparison proceeded in three stages of increasing rigour and the evolving result is itself informative.

{\bf Stage 1: Initial Comparison, 44-Gene Validation Set (Pre-RWR).} Both alternative methods significantly outperformed a naïve, expression-only global baseline (RRF: p = 0.017; LambdaMART: p = 0.008; paired bootstrap, 10,000 resamples). However, neither significantly outperformed our production system, which already incorporated mutation-impact scoring (RRF: p = 0.588; LambdaMART: p = 0.240).

{\bf Stage 2: Post-RWR Integration.} Having established that graph-topological information provides substantial, independently validated improvement (p = 0.0006), we re-ran the comparison to confirm this improvement was not simply a signal that RRF or LambdaMART could have exploited more effectively than our linear formula.

{\bf Stage 3: Fair Re-Comparison, All Methods Given Equal Feature Access.} We identified and closed a feature-parity gap. Copy-number and graph-centrality scores had not originally been exposed to RRF's fusion sources or LambdaMART's feature set. Once this was corrected, we repeated the significance test. Neither alternative method remained distinguishable from production (RRF: p = 0.270; LambdaMART: p = 0.995) and both scored numerically below production.

\begin{table}[t]
\centering
\begin{tabular}{|l|c|c|c|}
\hline
{\bf Stage} & {\bf RRF vs. Production} & {\bf LambdaMART vs. Production} & {\bf Outcome} \\
\hline
Stage 1 (Pre-RWR) & p = 0.588 & p = 0.240 & Not sig. \\
Stage 2 (Post-RWR) & p = 0.019* & p = 0.002* & Both sig. \\
Stage 3 (Fair comparison) & p = 0.270 & p = 0.995 & Not sig. \\
\hline
\end{tabular}
\caption{Significance testing across evaluation stages. * = significant at $\alpha = 0.05$. The trajectory shows that RRF and LambdaMART's apparent advantages at Stage 1 were attributable to production's prior omission of graph-topological evidence, not to inherent superiority of rank-fusion or tree-based learning-to-rank as combination mechanisms. Once all methods were given equal feature access, that advantage disappeared entirely.}
\label{tab:method-comparison}
\end{table}

The trajectory across these three stages is more informative than any single comparison in isolation. It reveals that the apparent advantage of RRF and LambdaMART at Stage 1 was specifically attributable to production's prior omission of graph-topological evidence, not to any inherent superiority of rank-fusion or tree-based learning-to-rank as combination mechanisms. Once production was given access to the same signal, that apparent advantage disappeared entirely. We retain the weighted-sum, gene-class-adaptive formula in production on two grounds: (i) no measurable performance disadvantage relative to more complex alternatives once all methods are compared fairly and (ii) substantially greater interpretability. Every individual score component in our formula maps directly to an auditable evidence source--a property neither RRF's fused rank nor LambdaMART's learned tree ensemble provides. We consider this interpretability essential for a tool intended to support defensible, reviewable experimental decisions in a regulated biopharmaceutical context.


\section{Graph-Based Ranking: Random Walk with Restart}

The project brief's first stated objective calls for the integration of multi-omic data into ``a unified, queryable knowledge graph to enable coherent cross-dataset analyses.'' Having established, in the preceding section, that pathway-membership coherence does not provide a useful ranking signal regardless of how it is computed, we turned to a structurally different question: rather than measuring whether a gene's pathway \textit{neighbours} are coherently expressed, could the \textit{topology} of the knowledge graph itself , the pattern of connections between genes, pathways, cell lines and mutations , provide information that our feature-based classical scorer could not?

{\bf Three Candidate Methods from the Network-Prioritisation Literature.}

We evaluated three established techniques for extracting a ranking signal from graph structure, each representing a distinct methodological family within the network-based gene- and entity-prioritisation literature.

Random Walk with Restart (RWR) \cite{kohler2008} simulates a random walker that begins at a seed node , in our application, the query gene  and at each step either follows a graph edge to a neighbouring node or returns to the seed with a fixed restart probability. The algorithm's steady-state visitation probability across every other node in the graph is interpreted as a measure of that node's structural proximity to the seed. We implemented RWR via personalised PageRank (NetworkX), a formulation that is mathematically equivalent to the canonical random-walk-with-restart algorithm.

Relational Graph Convolutional Networks (RGCNs) accommodate heterogeneous edge types (MEMBER\_OF, CONTAINS, EXPRESSED\_IN, HAS\_MUTATION) by learning edge-specific transformations. Although our lightweight, unsupervised RGCN trained on a link-prediction task failed to separate known-correct cell lines from candidates, we view this as a training-time constraint rather than a flaw in the RGCN architecture itself.

Metapath-guided walks constrain graph traversal to predefined, biologically meaningful path templates , for example, gene$\to$pathway$\to$gene$\to$cell line, or the direct gene$\to$mutation$\to$cell line path , rather than allowing an unguided walk to follow any available edge. We implemented a simplified, deterministic variant that explicitly enumerates and weights instances of each defined metapath, rather than the full stochastic training procedure of the original algorithm.

{\bf Standalone Evaluation: A Genuine but Insufficient Signal.}

Evaluated in isolation, using graph topology alone with no expression, protein, or mutation features whatsoever, RWR achieved a leave-one-out cross-validated Mean Reciprocal Rank of 0.1358 across our 44-gene validation set. This exceeded a na\"{i}ve, expression-only baseline (0.1096), demonstrating that the knowledge graph's structure genuinely encodes information relevant to cell-line suitability, independent of any direct genomic evidence. It did not, however, exceed our full classical scoring formula incorporating mutation-impact and copy-number evidence (0.1547 at that stage of development) , graph topology alone was informative, but less informative than direct biological evidence for the genes in our validation set.

This result alone was insufficient to justify replacing, or even meaningfully augmenting, the classical scoring approach. We report it, however, as an important intermediate finding: it establishes that the knowledge graph's contribution in what follows is not an artefact of double-counting information already present elsewhere in the pipeline, but a genuinely independent signal.

{\bf Integrated Evaluation: A Statistically Decisive Result.}

Rather than treating RWR as a candidate replacement for the classical formula, we integrated it as a \textit{complementary} evidence dimension, added to , not substituted for , the existing feature-based scoring. Class-specific weights were determined via grid search: 0.24 for tissue-specific genes, 0.28 for ubiquitous genes and 0.38 for loss-of-function genes, following the same weight-determination protocol used for every other scoring dimension in this dissertation.

The combined system achieved a leave-one-out cross-validated MRR of \textbf{0.2314}, compared to 0.1547 immediately prior to integration , an absolute improvement of 0.0767. A paired bootstrap significance test (10,000 resamples) confirmed this improvement was not attributable to chance: 95\% confidence interval [0.0280, 0.1338], \textbf{p = 0.0006}. This is the most statistically decisive result reported anywhere in this dissertation and the only comparison in our evaluation for which the confidence interval does not approach zero at either bound. (This figure reflects the system immediately following RWR integration, prior to the wild-type-preference correction described in the domain-expert validation section; the fully final production MRR, incorporating both corrections, is 0.2367.)

\begin{table}[t]
\centering
\small
\begin{tabular}{|l|c|c|}
\hline
{\bf Configuration} & {\bf LOO-CV MRR} & {\bf vs. pre-integration} \\
\hline
Classical formula alone (pre-RWR) & 0.1547 & , \\
RWR alone (no classical features) & 0.1358 & , \\
Classical formula + RWR (integrated)* & 0.2314 & +0.0767 (p = 0.0006) \\
\hline
\end{tabular}
\caption{Random Walk with Restart: standalone and integrated performance. * = Prior to the wild-type-preference correction; final production MRR is 0.2367.}
\label{tab:rwr-performance}
\end{table}

{\bf Interpreting the Result: Complementary, Not Redundant, Evidence.}

We interpret this finding as evidence that graph topology and direct genomic evidence (expression, mutation status, copy number) capture genuinely complementary, only partially correlated aspects of cell-line suitability. A cell line's position within the broader network of pathway and mutation relationships , which genes it is connected to and through what kind of relationship , provides signal that is structurally invisible to any of our direct evidence types considered individually and this signal proved most valuable precisely when combined with, rather than substituted for, that direct evidence.

A further test reinforced this interpretation. Having established that the classical-plus-RWR system significantly outperformed the classical formula alone, we asked whether this improvement could instead have been achieved more effectively by a more sophisticated combination algorithm , specifically, Reciprocal Rank Fusion and LambdaMART, described in the preceding section , given equal access to the same RWR signal. Once both alternative methods were given RWR as an available feature on equal footing with our production system, neither remained statistically distinguishable from it (RRF: p = 0.270; LambdaMART: p = 0.995). This confirms that the earlier apparent advantage of these alternative methods, observed before RWR integration, was specifically attributable to production's prior omission of graph-topological evidence , not to any inherent superiority of rank-fusion or tree-based learning-to-rank as a combination mechanism over our interpretable weighted-sum formula.

{\bf Relationship to the Project Brief.}

We regard this result as directly and substantively satisfying the brief's first objective. The knowledge graph in our final system is not a supplementary lookup structure consulted for auxiliary display information; it is a statistically validated, independently-contributing component of the core ranking mechanism, without which the system's overall accuracy would be measurably and significantly worse. We note this as a meaningful evolution over an earlier stage of development, at which the graph's only ranking contribution , pathway-membership coherence , had been found, through the investigation described in the preceding section, to provide no benefit and in several configurations active harm. The knowledge graph's eventual, validated contribution came not from the first mechanism attempted, but from correctly diagnosing why that first mechanism failed and identifying, from the established network-prioritisation literature, a structurally different way of extracting signal from the same underlying graph.

\section{Copy-Number Scoring for Amplification-Driven Oncogenes}

{\bf Motivation: A Third Mechanism of Gene Activation.}

Our gene-class-adaptive scoring framework distinguishes tissue-specific genes, best ranked by direct expression, from loss-of-function genes, best ranked by mutation status, on the grounds that different genes are selected as experimental models through different underlying biological mechanisms. In the course of this work, we identified a third such mechanism, distinct from both: a subset of oncogenes , \textit{MYCN}, \textit{ERBB2} and \textit{FGFR1} among the genes in our validation set , are canonically selected not on the basis of point mutation or elevated transcription alone, but of gene \textit{amplification}: the presence of many additional copies of the gene within the cell line's genome, driving protein overexpression through increased gene dosage rather than through altered protein sequence or transcriptional upregulation in the conventional sense.

This distinction matters in practice. A cell line with a single extra copy of \textit{MYCN} and one with a canonical, textbook-level amplification of one hundred or more copies may both register as ``expressing \textit{MYCN}'' under a conventional RNA-based expression score, yet only the latter is the kind of high-level amplicon a scientist would select as a \textit{MYCN}-amplified neuroblastoma model. Neither our expression-based scoring nor our mutation-impact scoring was designed to capture this distinction, since amplification is neither a change in transcript abundance per se nor a point mutation.

{\bf Data Source and a Release-Version Consideration.}

We sourced gene-level copy-number data from DepMap's public portal. Before selecting a specific file, we confirmed that our existing pipeline's other DepMap-derived files were drawn from two different quarterly releases, our expression data and ProfileID-to-ModelID mapping table from release 23Q4, but our somatic mutation calls from the more recent 24Q4, an inconsistency that had been present in the originally supplied project data and had gone unnoticed prior to this investigation. Because our identifier-resolution chain was built against the 23Q4 mapping table, we deliberately sourced the matching 23Q4 copy-number file rather than the more recent 24Q4 release, to avoid introducing a second release mismatch on top of the one already inherited in the mutation data.

{\bf A Scoring-Saturation Pitfall, Anticipated and Avoided.}

Having previously discovered and corrected a scoring-saturation defect in our mutation-impact scorer, in which a fixed ceiling caused biologically distinct cases (a borderline four-copy gain and a two-hundred-copy amplicon) to collapse to an identical maximal score, erasing the very distinction the score was intended to capture. We designed the copy-number scorer from the outset to avoid repeating this error. Rather than a score that saturates at a fixed log2 copy-ratio threshold, we implemented a continuous, slowly-saturating transformation in which increasingly extreme amplification continues to yield an increasing, if progressively smaller, score, ensuring that a modest gain and a definitive high-level amplicon remain numerically distinguishable rather than both mapping to the same maximal value.

We verified this design choice directly: across the neuroblastoma and breast-cancer lineages tested, no two cell lines with meaningfully different copy-number states received an identical score and the true extreme outlier in each lineage's amplicon distribution , a 295-copy \textit{MYCN} amplification in one neuroblastoma line and a 118-copy \textit{ERBB2} amplification in one breast-cancer line , correctly received the highest score within its respective lineage, distinct from and above the scores of more modestly amplified lines.

{\bf Integration and Validated Improvement.}

We restricted copy-number scoring to the three genes identified as amplification-driven in our validation set, applying it as an additional weighted evidence dimension (weight 0.30, determined by grid search) rather than to the tissue-specific gene class as a whole, since amplification is not the relevant selection mechanism for the majority of tissue-specific genes in our panel.

\begin{table}[t]
\centering
\small
\begin{tabular}{|l|c|c|}
\hline
{\bf Metric} & {\bf Before copy-number} & {\bf After copy-number} \\
\hline
Overall LOO-CV MRR & 0.2110 & 0.2142 \\
MYCN held-out rank & 0.250 & 0.333 \\
ERBB2 held-out rank & 0.010 & 0.023 \\
\hline
\end{tabular}
\caption{Effect of copy-number scoring on amplification-driven genes. These figures were measured at the point copy-number scoring was integrated, prior to the subsequent RWR and wild-type-preference additions; they isolate copy-number scoring's own contribution and are not directly comparable to the final production MRR of 0.2367.}
\label{tab:copy-number-improvement}
\end{table}

The improvement is modest in absolute terms, reflecting the small number of genes affected (three of forty-four in the final validation set), but consistent in direction and free of any measurable adverse effect on other gene classes, a pattern we regard as appropriately scoped: this scoring dimension was designed to correct a specific, identifiable mismatch between evidence type and biological mechanism for a small, well-defined subset of genes, not to serve as a general-purpose signal and its modest but real improvement for exactly the genes it targets, with no effect elsewhere, is the expected and desired outcome of a correctly-scoped addition.

{\bf A Recurring Methodological Pattern.}

We note that this is the second instance in this dissertation of a specific class of error, an unbounded or improperly-saturating score collapsing biologically meaningful differences into ties - being identified and corrected, first in the mutation-impact scorer and subsequently anticipated and avoided by design in the copy-number scorer. We regard this as evidence of a systematic risk in threshold- or ceiling-based scoring formulas applied to biological measurements with long-tailed distributions, of which raw mutation and copy-number data are typical examples and as justification for our general preference, established after the first occurrence, for continuous, slowly-saturating transformations wherever a hard biological ceiling does not genuinely exist.


\section{Investigated but Not Adopted: Pathway-Coherence Scoring}

Before arriving at the graph-based approach, we investigated whether pathway membership coherence--the extent to which genes sharing a KEGG pathway with the query gene are jointly expressed in a candidate cell line--could serve as a useful ranking signal. We tested four independent aggregation strategies for combining evidence across a gene's multiple pathway memberships. All four consistently failed to improve ranking accuracy for tissue-specific and ubiquitous genes and in several configurations measurably degraded it. We report this investigation in full because we regard the eventual mechanistic explanation for this consistent failure and its direct connection to the mutation-impact scoring approach that ultimately succeeded, as a genuine and valuable scientific finding rather than a gap in the work.

{\bf Strategy 1: Naïve Pathway Co-Expression Count.} Our initial approach measured, for a given target gene and candidate cell line, the proportion of the gene's KEGG pathway co-members that were also expressed above a fixed threshold in that cell line. The intuition: a cell line with coherently active pathway is a stronger model than one where only the gene itself is expressed. Tested via exhaustive grid search across pathway weight values from 0.00 to 0.50, this measure showed clean, monotonically decreasing relationship with ranking accuracy for tissue-specific genes. Every weight increment above zero reduced Mean Reciprocal Rank, without exception across the tested range.

{\bf Strategy 2a: ssGSEA Union Across All Pathways.} We hypothesised naïve count's failure arose from coarse, threshold-based treatment of expression. We implemented single-sample Gene Set Enrichment Analysis (ssGSEA) \cite{barbie2009ssgsea}, measuring whether a pathway's member genes are collectively enriched toward the top of a cell line's full ranked transcriptome.

For EGFR, which is a member of twenty-two KEGG pathways, we tested combining evidence as a union: one set containing all co-members from all twenty-two pathways, approximately 1,900 genes total. For canonical EGFR-mutant reference lines HCC827 and NCI-H3255, this placed them at 41st and 42nd percentile respectively--better than chance, but far from confident top-ranked results. Dilution effect dominates. A gene set of nearly two thousand members, drawn predominantly from large, biologically generic pathways such as MAPK signalling and PI3K-Akt signalling (each exceeding 200 genes), is dominated numerically by genes with only distant relationship to the query gene's specific biology.

{\bf Strategy 2b: Maximum ssGSEA Across Independently-Scored Pathways.} To avoid dilution, we scored each pathway membership independently and took the maximum resulting score across all pathways, reasoning that a cell line coherent with even one relevant pathway should be considered a good candidate. This performed worse than the union strategy and revealed a specific, instructive statistical error. Taking the maximum across twenty-two independently-distributed pathway scores per cell line constitutes an uncontrolled multiple-comparisons procedure. Under the null hypothesis that no pathway is genuinely informative, the expected value of the maximum of twenty-two independent uniform draws already approaches 0.96, meaning nearly every cell line appears to score highly on some pathway purely by chance.

We confirmed this diagnosis in the actual output: eight distinct cell lines tied at perfect maximum score of 1.000, each achieving that via a different pathway. For three of the four known EGFR-mutant validation lines tested, the pathway responsible for the winning score was not ErbB signalling--the pathway containing EGFR itself and the one domain experts would consider most relevant--but an unrelated pathway such as GnRH signalling, adherens junction, or endocytosis. This strongly indicates statistical noise rather than coherent biological signal.

{\bf Strategy 2c: Specificity-Weighted Combination.} As a compromise, we weighted each pathway's contribution inversely to its size before summing them, reasoning that small, specific pathways (ErbB signalling, 86 genes) should contribute more than large, generic ones (PI3K-Akt, 350+ genes). Even with inverse-size weighting, EGFR's most specific and arguably most relevant pathway (ErbB signalling) received only approximately 8\% of the total combination weight, with the remaining 21 pathways accounting for 92\%. A weighted sum over 22 terms, however weighted, remains fundamentally an averaging operation and the specificity weighting proved too weak relative to the sheer number of competing, less-relevant terms.

{\bf Strategy 2d: Single Canonical Pathway (Diagnostic Only).} As a final diagnostic, we tested scoring EGFR using only its single most specific and biologically canonical pathway, ErbB signalling, in isolation. This was, by substantial margin, the best-performing configuration tested: known EGFR-mutant reference lines HCC827 and NCI-H3255 rose to 92nd and 81st percentile respectively. We do not propose this as production method, as it requires identifying a single canonical pathway for arbitrary genes--straightforward for EGFR, considerably less so for genes without one obviously dominant pathway membership. We report it as diagnostic result, confirming that failure of Strategies 2a–2c lies in the combination mechanism across multiple pathways, not in pathway membership being inherently uninformative.

\begin{table}[t]
\centering
\small
\begin{tabular}{|l|c|c|c|}
\hline
{\bf Strategy} & {\bf HCC827 Percentile} & {\bf NCI-H3255 Percentile} & {\bf Problem} \\
\hline
Strategy 2a (Union) & 41 & 42 & Dilution from large pathways \\
Strategy 2b (Maximum) & 55 & 27 & Multiple-comparisons inflation \\
Strategy 2c (Specificity-weighted) & 43 & 44 & Still diluted despite weighting \\
Strategy 2d (Canonical only) & 92 & 81 & Best result, but unmasterable \\
\hline
\end{tabular}
\caption{Pathway-combination strategy results for EGFR against known EGFR-mutant reference lines. All four strategies failed to produce confident top-ranked results. Strategy 2d diagnostic result shows that the failure is not inherent to pathway evidence, but to combination mechanisms that dilute or average over multiple pathway memberships.}
\label{tab:pathway-strategies}
\end{table}

{\bf Mechanistic Explanation.} Taken together, these four results point to a single underlying explanation. The receptor tyrosine kinase and oncogene-class genes represented in the validation set--EGFR foremost among them--are canonically selected as experimental models because an activating point mutation renders the encoded protein constitutively active through a post-translational, phosphorylation-cascade signalling mechanism. This mode of pathway activation does not require and in general does not produce, coordinated transcriptional upregulation of the pathway's downstream gene set. Many kinases and adaptor proteins in a signalling cascade are constitutively expressed at baseline levels in most cells, with their activity--not their expression level--determined by the mutation status.

Transcriptomic pathway-enrichment methods, including ssGSEA, are structurally unable to detect a mode of biological activation that manifests post-translationally rather than transcriptionally. This fundamental limitation explains the uniform failure of Strategies 1 and 2a–2c across every combination mechanism tested, independent of the specific aggregation strategy employed.

{\bf Motivated Alternative Approach.} This mechanistic understanding directly motivated the subsequent and considerably more successful, adoption of mutation-impact scoring (described in Chapter 3). Rather than attempting to infer pathway activation from a downstream transcriptional signature that the relevant class of genes does not reliably produce, mutation-impact scoring detects the causal genomic lesion directly. It identifies that HCC827 carries EGFR exon 19 deletion and assigns this a high-impact score, capturing the constitutively active kinase regardless of transcriptional consequences.

{\bf Scientific Value of Negative Results.} We regard the connection between this negative result and the positive alternative as itself a meaningful contribution to the broader scientific literature. It demonstrates that understanding why an evidence type is uninformative for a given biological mechanism enables identification of the correct alternative evidence type, rather than arriving at it by trial and error. For researchers considering similar pathway-integration approaches for other genes or diseases, this work provides evidence that transcriptomic enrichment cannot reliably detect post-translational activation mechanisms, clarifying the boundary of applicability for these methods and potentially saving future effort on similarly-motivated but ultimately ineffective strategies.

\section{Copy-Number and Amplification-Driven Genes}

A distinct subset of oncogenes--MYCN, ERBB2, FGFR1 and others in the validation set--are canonically selected as experimental models not on the basis of point mutation, but of gene amplification: the presence of many extra copies of the gene, driving overexpression through increased gene dosage rather than altered protein sequence. To address this mechanism, distinct from both expression-driven and mutation-driven scoring already described, we developed a copy-number scorer using DepMap gene-level copy-ratio data (release 23Q4).

This scorer applied a continuous, slowly-saturating transformation calibrated such that a modest four-copy gain and a high-level, hundred-copy amplicon remain distinguishable rather than collapsing to identical maximum score--a saturation failure mode we had previously diagnosed and corrected in mutation-impact scoring and deliberately avoided in this later addition to prevent repeating the same error.

This scoring dimension, weighted at 0.30 for the three amplification-driven genes identified in the validation set, improved their individual held-out reciprocal ranks and contributed a measurable, if modest, improvement to overall cross-validated MRR.

\section{Validation Against Independent Domain-Expert Expectations}

Statistical validation characterises system behaviour against the team's own curated validation set. This form of validation, while methodologically rigorous, shares a structural limitation: the validation set is constructed by the same team that built the system and cannot by its nature rule out the possibility that both share a common blind spot. To address this, we sought an independent source of ground truth: unprompted, domain-expert predictions supplied directly by an AstraZeneca scientist unconnected to system development.

These predictions took the form of a short list of genes, each accompanied by stated expectations of how a well-behaved cell-line recommendation system ought to treat it, given known biological properties of that gene.

{\bf The Housekeeping-Gene Prediction: GAPDH.} The domain expert noted that GAPDH is expressed at broadly similar levels across nearly all human cell lines and predicted that a correctly-behaving system should therefore show little differentiation in scores for this gene. This is a genuinely falsifiable prediction, one that a system reverse-engineered to ``look confident'' on every query would be liable to fail.

We tested it directly by examining the full score distribution--not merely the top-ranked results--returned for GAPDH across all candidate cell lines. The distribution was markedly flatter than that of any tissue-specific gene tested. The gap between the first and fiftieth-ranked score for GAPDH was the smallest observed across all comparator genes, by a margin exceeding a factor of two. We interpret this as direct confirmation that the system does not artificially manufacture differentiation where biology does not support it--an important property for a tool intended to support, rather than substitute for, expert judgement. A system appearing confident regardless of underlying evidence would actively mislead.

{\bf The Sparse-Evidence Prediction: CD3E.} A second prediction concerned CD3E, a well-established T-cell marker in published literature, for which specific integrated data sources (proteomics and GEO coverage in particular) were expected to be sparse. This tested a different and harder property: sparse data coverage should be correctly distinguished from biological unimportance or poor expression. A naïve system would conflate ``we have little data'' with ``there is little signal.''

We queried the system for CD3E against Jurkat, a canonical T-cell line. The system correctly identified and reported that CCLE proteomics data was absent for this gene–cell-line pair, explicitly naming this as a data limitation. Simultaneously, it reported strong corroborating evidence from available sources (92nd and 96th percentile RNA expression from HPA and DepMap respectively) and recommended Jurkat positively on that basis. The generated justification did not conflate absence of proteomic confirmation with absence of biological relevance; it named the missing evidence type explicitly while surfacing strong available evidence.

{\bf A Genuine Defect Surfaced: TP53 Wild-Type Preference.} Extending this face-validity exercise to genes beyond those explicitly named by the domain expert, we queried TP53 in a colorectal-cancer context and manually inspected top-ranked results against domain knowledge. HCT116, the field's standard TP53 wild-type reference line, used throughout cancer biology literature specifically because it retains functional, unmutated p53, was ranked below two TP53-mutant lines, including one carrying the canonical hotspot mutation p.R273H.

This was a genuine defect, not a modelling limitation. Our scoring formula applied no penalty for mutation status when scoring ubiquitously-expressed genes. A damaging TP53 mutation contributed no negative signal despite the gene's canonical use case in this context being selection of an unmutated control line. We diagnosed the scope precisely before correcting. Querying our mutation dataset directly: TP53 carries damaging variant in 1,040 cell lines in our data, compared to 6–40 for every other gene in the same ubiquitous category--a 20- to 170-fold outlier that confirmed this was TP53-specific rather than a defect affecting the ubiquitous gene class as a whole.



We introduced a small, fixed wild-type-preference penalty (-0.05), applied only to TP53 and re-tested.

\begin{table}[t]
\centering
\begin{tabular}{|l|c|c|}
\hline
{\bf Rank} & {\bf Before Correction} & {\bf After Correction} \\
\hline
1 & SW948 (TP53-mutant) & SW48 (TP53 wild-type) \\
2 & SW48 (TP53 wild-type) & SW948 (TP53-mutant) \\
3 & SW620 (TP53-mutant, p.R273H) & SK-CO-1 (TP53 wild-type) \\
HCT116 Rank (of 75 total) & 32 & 22 \\
\hline
\end{tabular}
\caption{Top-3 results for TP53 in colorectal-cancer context, before and after wild-type-preference correction. Held-out cross-validated MRR for ubiquitous gene class improved from 0.0633 to 0.1216. TP53's own held-out reciprocal rank improved from 0.100 to 0.333. Other ubiquitous-class genes (PARP1, CDK4, CCND1) showed exact-zero change, confirming the correction was properly scoped to the specific gene responsible for the anomaly.}
\label{tab:tp53-correction}
\end{table}

We confirmed the fix produced no adverse effects: ubiquitous-class held-out MRR improved from 0.0633 to 0.1216 and TP53's own held-out reciprocal rank improved from 0.100 to 0.333. The three other ubiquitous-class genes (PARP1, CDK4, CCND1) showed exact-zero change, confirming the correction was correctly scoped to the specific gene responsible rather than being either too broad or too narrow.

{\bf Methodological Significance.} We consider this sequence of events--an independently-sourced, falsifiable prediction leading directly to discovery and correction of a genuine defect that our own curated validation set did not expose to be among the more methodologically significant findings in this dissertation, independent of its effect on any single numerical result. It demonstrates a general principle applicable beyond this specific project: a validation set constructed by the same team that built the system cannot, by its nature, rule out shared blind spots and qualitative, domain-expert-informed spot-checking is not merely supplementary but a necessary complement to quantitative cross-validation, capable of surfacing an entire category of error that aggregate statistical metrics are structurally unable to detect.

\section{Summary}

Through systematic testing, critical bug correction, rigorous comparison against alternative methods and independent domain-expert validation, we have demonstrated that CellSelector Omics produces reliable, interpretable rankings grounded in evidence. The gene-class hypothesis that different genes require fundamentally different ranking logics is supported by convergent evidence from learned weights, cross-validation performance, mechanistic understanding and domain-expert validation. The investigation into pathway-coherence scoring, though yielding no improvement, provided mechanistic insight that directly motivated the successful mutation-impact approach.


\chapter{Conclusion}
\label{chap:conclusion}

\section{Achievement Against Project Aims}

This project replaced the manual, expertise-dependent and error-prone process of cell-line selection with an automated system that is evidence-backed, biologically-aware and interpretable to domain scientists. All five aims stated in Chapter 1 were met and substantially exceeded in scope, depth and validation rigor.

{\bf Aim 1: Integrate Multiple Public Multi-Omics Datasets.} Achieved and Exceeded. Five databases (HPA, DepMap, GEO, CCLE, KEGG) have been harmonized against the canonical Cellosaurus identifiers, encompassing 2,076 cell lines and 222 million expression values. It is a heroic level of integration – the data store covers more than 200 million data points, but still keeps the full granularity. And it is not easy technical task. There are incompatible units (raw counts, TPM, log-TPM), identifiers (gene symbols, Ensembl IDs, RefSeq) and different coverage patterns in the used data sources. The naive approach would require full coverage from all sources and dropping of 60-70\% of cell lines. However, our approach makes it possible to preserve partially covered cell lines by averaging the scores across sources rather than requiring full coverage. This is exactly how our approach tackles one of the problems of the current approaches. The unified data store is created based upon Parquet files for fast querying of the entire 222-million-size dataset using only laptop without database server. Access is crucial for adoption and reproducibility.

{\bf Aim 2: Enable Gene-Presence and Gene-Exclusion Ranking.} Achieved with a user-friendly interface. Gene-inclusion queries provide single-gene and multi-gene searches with principled fusion of evidence by returning ranking cell lines arranged by evidence for target gene expression. Because scientists eliminate cell lines based on restrictions rather than preferences, gene-exclusion queries are implemented as hard filtering rather than soft penalties. This design decision reflects the use case. Regardless of ranking score, a cell line with high expression of the melanoma marker is unsuitable for lung cancer research. Each of the four separate criteria supported by the exclusion interface—microsatellite instability, chromosomal instability, metabolite levels and miRNA expression—is user-configurable rather than hardcoded. Instead than being constrained by default settings, this flexibility allows scientists to define exclusion levels depending on their unique experimental design. Complex queries like "include gene A," "exclude genes B and C," "require mutation in D," and "exclude cell lines with MSI" are all supported by boolean logic. The technology advances from point-and-click search to a query language for cell-line biology because to this expressiveness.

{\bf Aim 3: Quantify Confidence Despite Sparse, Heterogeneous Data.} Achieved through explicit missing-data handling and three validations.Uncertainty is revealed rather than concealed by a gene-class aware scoring model with explicit missing-data handling. Significantly, this tracking of uncertainty turned out to be helpful during development. Each of the three different big problems that were found and fixed would have reduced ranking accuracy if they had not been disclosed. 94\% of HPA mappings and 80\% of GEO mappings were incorrect due to Bug 1 (HPA/GEO identification misattribution), a silent problem that would have been missed without close inspection. Because of a predetermined iteration sequence, Bug 2 (non-deterministic tie-breaking) broke reproducibility criteria by causing the identical query to get different replies on successive executions. When mutation-impact scores hit 1.0, Bug 3 (score-saturation artifact) artificially increased confidence in loss-of-function genes, obscuring weaker evidence from other sources. Every bug was found through -validation and careful testing. This shows how important thorough validation workflows are, for every bug. In order to provide users with an informed assessment of the results, the system now records confidence independently of ranking. The weight of a cell line ranked first based solely on RNA evidence differs from that of a cell line ranked first with high confidence based on four independent data sources.

{\bf Aim 4: Provide Transparent, Evidence-Cited Justifications.} Achieved through agentic RAG with strict sourcing constraints. To find pertinent evidence, an agentic RAG pipeline searches PubMed, KEGG, Cellosaurus and the underlying omics data. Explicit instructions to synthesize reasons citing solely from these retrieved sources—never creating proof or speculating—place restrictions on the language model. This design distinguishes between explanation (which is flexible in narrative but based on recovered evidence) and ranking (which is deterministic, explainable and auditable). The design of the constraints is crucial. A basic method would enable the language model to generate assertions that sound credible but cannot be independently verified. Rather, each statement in the reasoning must link to a specific data point from the omics layers, a retrieved PubMed abstract, a KEGG pathway entry, or a Cellosaurus record. In stakes biomedical settings, where false statements, about cell lines could ruin research this design avoids hallucinations. The system mixes domain- language well and keeps correctness. The design is clear cell lines are accurate and the system stays reliable.

{\bf Aim 5: Deliver an API and User Interface with Exportable Reports.} Production‑grade deployment has been achieved. Programmatic integration into workflows is enabled by a FastAPI backend that offers a stateless JSON API. The API is organized into four routers: genes, ranking, cell lines and filters. Interactive exploration is possible with a React web frontend. The frontend includes SearchBar, FilterPanel, ResultsTable and DetailPanel components. Evidence visibility is a design choice. Every ranked cell line is shown together with an evidence breakdown by source. This lets scientists decide if data sources support or contradict a recommendation. Results can be exported in PDF, CSV or JSON format. Reports that contain figures, tables and evidence citations, for publishing are created by using the PDF export. Continuous deployment is secured by GitHub Actions CI/CD. After code is merged the new version is. Deployed within two or three minutes. The deployment runs on GCP server. This setup delivers the security and availability that a shared research infrastructure requires. Sub-200 millisecond load time ensures interactive response despite multi-gene queries.

\section{Key Scientific Finding: The Gene-Class Hypothesis}

The most important technical discovery is that ranking accuracy is determined by gene class rather than a single global weighting. This realization, which represents a mechanistic understanding of the ranking problem itself, came about as a result of research into systematic failure modes.

{\bf Discovery Through Failure.} Tissue-specific genes were appropriately ranked by an early model that averaged expression evidence evenly across all genes, while findings for ubiquitous and loss-of-function genes were biologically absurd. The technique identified high-expressing lines as the best options by ranking cell lines according to their absolute expression level when searching for common genes like GAPDH or ACTB. However, almost all cell lines express GAPDH at about identical levels; absolute expression does not offer a discriminative signal for cell line selection. A cell line ranked first for GAPDH was no better than one ranked 50th; both express it at housekeeping levels. When queried for loss-of-function genes like TP53, the system identified high-expressing lines, which is precisely wrong. Wild-type TP53 is expressed in nearly all cell lines. Cell lines with high TP53 expression are not appropriate for studying tumour suppressor loss. Cell lines with mutant, non-functional TP53 are the scientifically relevant models. These systematic failures revealed that a single weighting scheme cannot work across all genes; different genes require different evidence hierarchies.

{\bf Three-Way Classification.} Investigation of failure modes led to classification into three biologically meaningful categories. Tissue-specific genes like EGFR or ERBB2 are studied because of their high, cancer-specific expression. They are implicated in cancer development through gain-of-function: overexpression drives proliferation. For these genes, ranking by expression level is correct. Ubiquitous genes like GAPDH or ACTB are expressed at similar levels across all cell lines and are not targets of cancer-specific study. They appear in queries as negative controls or experimental normalisation targets. Querying for ubiquitous genes should return cell lines suitable for the surrounding experimental context, not cell lines high in the ubiquitous gene itself.

{\bf Evidence From Learned Weights.} The strength of evidence for this classification lies in convergence of multiple independent lines of evidence. Most compellingly, the learned per-class weights recover biological expectations despite the optimizer having no biological instruction. The optimizer was constrained only to maximise Mean Reciprocal Rank on the validation set, with no guidance about gene classes or biological reasoning. Yet it converged on markedly different weightings:

Tissue-specific genes: RNA weight 0.60, Protein weight 0.10, Quality weight 0.20, Mutation-impact weight 0.10. The optimizer assigned nearly all weight to RNA expression, the relevant signal for tissue-specific genes.

Ubiquitous genes: RNA weight 0.50, Protein weight 0.05, Quality weight 0.35, Mutation-impact weight 0.10. Expression plays a smaller role; quality and consistency of evidence carry greater weight, reflecting that ranking by absolute expression is unreliable.

Loss-of-function genes: RNA weight 0.003, Protein weight 0.048, Quality weight 0.046, Mutation-impact weight 0.80. A 270-fold difference between expression and mutation-impact weighting. The optimizer independently discovered that expression is nearly irrelevant for loss-of-function genes; mutation status dominates. This is precisely what a biologist would predict. The emergence of this insight from pure optimisation on ranking performance, without biological instruction, provides the strongest evidence that the gene-class distinction reflects a fundamental property of the ranking problem rather than a convenient categorisation.

{\bf Magnitude and Stability of Improvement.} The order-of-magnitude differences in per-class Mean Reciprocal Rank reflect fundamental problem properties, not method deficiency. Ubiquitous genes achieved MRR 0.0062 under global weighting but 0.1216 under class-aware weighting, a 19.6-fold improvement. Loss-of-function genes achieved 0.0041 globally but 0.1968 with class weighting, a 48-fold improvement. Tissue-specific genes achieved 0.1547 globally but 0.3300 with class weighting, a 2.1-fold improvement. These ratios match what we know about biology. They can be repeated in runs and stay steady across all cross‑validation folds. They show progress in making ranking match biology. They are not noise or small‑sample quirks.

{\bf Comparison Against Alternatives.} When I compared the statistical model to many machine‑learning methods I saw that the traditional model is easier to understand and still holds its own against the more advanced methods if we use the same features. I tested Reciprocal Rank Fusion and LambdaMART an algorithm, with the identical feature set. Stage 1 gave a p‑value of 0.588 Stage 3 gave 0.270 so there is no difference, between RRF and the traditional model. Stage 1 p=0.240 and Stage 3 p=0.995 indicate no significant difference between LambdaMART and the conventional model. This result has two implications. First, the resilience of the classical model is not improved by more complex methods. Second, the primary sources of progress are gene-class distinction and feature engineering rather than computational sophistication. An easier-to-understand model that outperforms black-box alternatives is necessary for adoption in regulated research environments.

{\bf Domain-Expert Validation.} Validation against domain-expert forecasts revealed a critical biological error that was missed by quantitative metrics. The system ranked wild-type TP53 cell lines above TP53- lines. High baseline expression of genes allowed wild-type TP53 to dominate rankings. This was not appropriate for tumor suppressor studies. Implementing a instance-level penalty of -0.05 for wild-type TP53 corrected this inversion. This change doubled the ubiquitous-class MRR from 0.0633 to 0.1216. This shows that face-validity checking by domain experts is essential. It is essential alongside evaluation. It helps to catch domain-level flaws that metrics overlook. This shows that face-validity checking by domain experts is essential. It is essential, alongside evaluation. It helps to catch domain-level flaws that metrics overlook.

\section{Negative Results as Scientific Contribution}

A systematic investigation into pathway-coherence scoring tested four independent methods to combine pathway activity with gene-expression ranking. Ranking accuracy was not improved over the baseline by any of the four approaches. When looked at from a point of view these unfavorable outcomes become a real scientific contribution.

{\bf Strategy 1: Naïve Co-Expression Count.} In a cell line the first method counted how many genes, in a pathway were expressed with the query gene. If many ErbB signaling genes are expressed similarly to EGFR, it is likely that the cell line has active ErbB signaling. In practice, this method produced MRR that dropped monotonically as weight increased. Every rise above zero was harmful. The MRR curve continuously declined, suggesting that pathway co-expression is actively misleading rather than just uninformative. The root cause is often cell-line-specific and noisy transcriptional co-regulation of genes in the same pathway. Membership in a pathway does not require coordinated expression. It's a fairly simple method.

{\bf Strategy 2a: ssGSEA Union Across Pathways.} The second method treated the pathway of the query gene as the target set and using single-sample gene set enrichment analysis (ssGSEA) to assess pathway activity within each cell line. For EGFR (ErbB signalling pathway member), the system computed ssGSEA scores for the entire 155-gene ErbB pathway within each cell line, ranked cell lines by ErbB pathway score. The problem: ssGSEA acts as a regulariser across pathway members. A cell line with specific EGFR upregulation but moderate expression of other pathway members receives a moderate pathway score, diluted by the majority. The top-ranked cell lines for EGFR query positioned at 41st and 42nd percentile in ErbB pathway enrichment, far below cell lines with generic pathway activity. The dilution effect dominates; large pathway sets (1,900 genes in KEGG pathways) cannot detect specific signal from individual genes.

{\bf Strategy 2b: Maximum ssGSEA Across Pathways.} To address dilution, the third approach took the maximum ssGSEA score a cell line achieved across any pathway, on the theory that activation of any upstream pathway might facilitate the query gene's function. As a result, multiple-comparisons inflation emerged as a distinct issue. Cell lines that achieved high ssGSEA in pathways unrelated to the query were given high scores. Eight cell lines tied with a perfect 1.0 score for the EGFR inquiry by activating completely separate pathways (endocytosis, adherens junction and GnRH signaling). These pathways have no biological connection to EGFR and the high scores are not indicative of a true signal but rather of noise and multiple comparison effects. The approach is not trustworthy.

{\bf Strategy 2c: Specificity-Weighted Combination.} The fourth method gave ErbB signaling a higher weight than GnRH signaling in an effort to weight pathways according to their relevance to the query gene. In actuality, the optimizer allocated weights so that just 8\% of the total weight went to ErbB signaling. Other routes predominated. The explanation is that whereas generic pathways like cell-cycle or metabolism exhibit high ssGSEA scores in the majority of cell lines, ErbB signaling is particular (few cell lines display strong signal). Frequent, high-signal paths are preferred over particular ones via optimization. This bias was not corrected by specificity weighting.

{\bf Strategy 2d: Single Canonical Pathway.} The final method employed only the canonical pathway for the query gene (ErbB signaling for EGFR) and did away with cross-pathway comparison. This removed multiple-comparison noise but exposed a different problem. Target cell lines (HCC827, NCI-H3255, both known EGFR-addicted) ranked at 92nd and 81st percentile of ErbB pathway enrichment. The problem: cell lines can have high EGFR expression independent of whole-pathway activation. HCC827 expresses EGFR because of an activating mutation (exon 19 deletion) that drives autocrine signalling; other pathway members are not necessarily coordinately upregulated. Pathway enrichment misses this mutation-driven activation.

{\bf Mechanistic Insight: Post-Translational Activation.} The systematic failure of all four pathway methods points to a fundamental mechanistic insight. Post-translational phosphorylation-cascade activation does not produce coordinated transcriptional upregulation. In EGFR-addicted cell lines, constitutive EGFR signalling occurs through receptor autophosphorylation and downstream kinase activation. These events are post-translational.

{\bf Motivated Alternative Approach.} This mechanistic understanding directly motivated the successful alternative: mutation-impact scoring. Rather than attempting to infer pathway activation from transcriptomics, the system directly captures the effect of gain-of-function mutations (EGFR exon 19 deletion, EGFR L858R, ERBB2 amplification) that activate kinase cascades independent of transcription. Mutation-impact scoring achieved substantially better performance than all pathway methods. The negative result, understood mechanistically, contributed directly to the project's success.

{\bf Scientific Value of Negative Results.} Negative results understood mechanistically represent genuine scientific contribution. The project eliminated an intuitive but ineffective approach, clarified the boundary of applicability for transcriptomic pathway methods and directly motivated successful alternative. For researchers considering similar pathway-integration approaches: this work provides evidence that transcriptomic enrichment cannot reliably detect post-translational activation mechanisms, saving future effort on similarly-motivated strategies. The mechanistic explanation (transcriptomics invisible to phosphorylation cascades) generalises beyond this specific project.

\section{Limitations and Future Work}

The project has explicit limitations that should inform future use and development. These are not weaknesses but rather clear boundaries that guide next steps.

{\bf Experimental Validation Absent.} The system's predictions have never been experimentally validated. No cell lines selected by the system were cultured, subjected to the planned experiment and independently assessed for biological appropriateness. This is the highest-priority limitation and the most important future work. Experimental validation would answer the critical question: does ranking accuracy (quantitative, internal, measured against Cellosaurus annotations) translate to biological utility (qualitative, external, measured by experimental success)?

Expected outcomes: if precision $>80\%$ and recall $>70\%$, the system is ready for adoption. Ranking is inconsistent and has to be debugged if precision $<50\%$. The method identifies some good lines but misses others if recall is low while precision is high; expanding the data sources might be beneficial. Certain gene classes require more research if class-specific performance differs.

{\bf Validation Set Small.} Despite being carefully chosen for diversity, the 44-gene evaluation set is modest in comparison to the genome scale (20,000+ protein-coding genes). Particularly for rare gene functions or genes with low data coverage, improvements in ranking accuracy may not translate to novel genes outside the validation set. Confidence in generalization would be strengthened by cross-validation using other gene sets.

Future research should: (1) acquire independent evaluation sets from partners (Novartis, GSK) with internal expertise in cell-line selection; (2) evaluate class-specific generalization through stratified cross-validation by gene ontology category (kinases, transcription factors, cell-surface receptors); and (3) compare system rankings with past cell-line selections at partner organizations to determine whether the system would have suggested the lines that experts selected.

{\bf Coverage Constraints Limit Scope.} Proteomics coverage is sparse (18.1\% of 2,076 cell lines have CCLE proteomics data, n=375). Mutation data is incomplete (60.1\% have annotated mutations, n=1,246). For cell lines lacking proteomics or mutation data, recommendations rest entirely on RNA evidence. These gaps mean recommendations are strongest for lines with high coverage and weaker for lines with sparse data. Scientists should be aware of coverage when interpreting results.

\begin{table}[h]
\centering
\begin{tabular}{|l|c|c|l|}
\hline
{\bf Data Source} & {\bf Cell Lines} & {\bf Coverage} & {\bf Implication} \\
\hline
DepMap RNA-seq & 1,399/2,076 & 67.4\% & Primary expression layer \\
HPA RNA-seq & 1,198/2,076 & 57.8\% & Independent RNA validation \\
CCLE proteomics & 375/2,076 & 18.1\% & Critical bottleneck \\
DepMap mutations & 1,246/2,076 & 60.1\% & Adequate for most queries \\
\hline
\end{tabular}
\caption{Coverage of key data layers. Transcriptomics near-complete (67-58\%); 
proteomics is critical constraint (18.1\%); mutations adequate (60.1\%). Expanding 
proteomics coverage would yield highest marginal return.}
\label{tab:coverage}
\end{table}

These coverage gaps mean recommendations are strongest for lines with high multi-modal 
coverage and weaker for lines with sparse data. Scientists should be aware of coverage 
when interpreting results.


Future work should: (1) integrate OncoKB mutation annotations, adding clinical interpretation of mutations; (2) integrate COSMIC mutation data, covering additional cell-line variants; (3) expand proteomics coverage by including independent proteomics studies beyond CCLE; (4) develop probabilistic confidence intervals that expand when coverage is sparse, making uncertainty explicit to users.

{\bf Scope Limited to Cancer Cell Lines.} All 2,076 cell lines in the system are cancer-derived. The system's applicability to primary cell cultures, patient-derived xenografts, organoids or non-cancer cells (primary hepatocytes, neural cells, immune cells) is completely unexplored. Cancer biology has distinctive features (aneuploidy, oncogenic mutations, altered metabolism) that might make recommendations inapplicable to normal cells. Different rules might apply for selecting normal cell lines.

Future work: (1) obtain primary-cell RNA-seq data from databases like GTexPortal and HCA; (2) retrain gene-class weights on primary cells; (3) develop separate scoring for cancer vs. primary distinction; (4) validate on organoid and PDX models, which are intermediate between cancer lines and primary cells.

{\bf Tissue-Specific Reliability Uncertified.} The system was optimised on the mixed population of all 2,076 cell lines. Different tissues (lung, breast, ovary, colorectal) have different mutation landscapes and expression patterns. Ranking reliability might vary by tissue. A tissue-specific optimisation of gene-class weights might improve performance for tissue-specific queries.

Future work: (1) stratify validation set by tissue of origin; (2) refit gene-class weights separately for lung, breast, ovary, colorectal (the four most common); (3) measure whether tissue-specific weights improve within-tissue ranking while degrading cross-tissue ranking; (4) expose tissue choice to users, allowing tissue-specific optimisation when appropriate.

{\bf Default Weights Unvalidated on External Data.} The per-class weights (tissue-specific: RNA 0.60, ubiquitous: RNA 0.50, loss-of-function: RNA 0.003) and RNA/protein balance (0.7/0.3) were learned on the internal 44-gene validation set. Generalisation to other gene sets is assumed but unverified. Different research contexts might require different weightings.

Future work: (1) perform sensitivity analysis: how much does ranking accuracy change if each weight is perturbed $\pm 20\%$?; (2) elicit weights from domain experts (10-15 experienced cell biologists); (3) compare expert weights against learned weights; (4) provide users with weight-adjustment interface, allowing labs to optimise for their specific use cases.

\subsection*{Synthesis: Practical Guidance for System Use}

{\bf High-confidence recommendations:} Tissue-specific genes (EGFR, HER2, PD-L1) 
in well-characterised cancer types (lung, breast, ovary) where DepMap RNA + HPA 
RNA + CCLE proteomics all available. Tissue-specific MRR (0.3300) and P@10 (0.48) 
both support confidence.

{\bf Medium-confidence recommendations:} Loss-of-function genes (TP53, BRCA1, PTEN) 
where mutation data available. Gene-class weighting recovers biological principles, 
but recommendations warrant cross-checking against domain knowledge.

{\bf Low-confidence recommendations:} (1) Ubiquitous genes; (2) sparse data coverage; 
(3) rare cancer subtypes; (4) non-cancer cells. Treat system as starting point, not 
validated shortlist. Pair with expert review or pilot experiments.

{\bf Highest-value use:} System most valuable for tissue-specific genes in common 
cancers, automating tedious manual work across five databases. Its value is in 
consistency and speed, not in discovering novel rankings.


{\bf Prioritised Future Work Sequence.}

Priority 1: Experimental validation. This is the highest-value activity. It addresses the essential query, "Do predictions work in practice?" and offers useful input to direct system enhancements. Set aside one or two experimental scientists and three to four months. Anticipated result: certification of the system's adoption readiness or discovery of certain flaws that need to be fixed.

Priority 2: Using OncoKB and COSMIC to expand mutation-data. The most significant signal for loss-of-function genes is mutation data (weight 0.80). Loss-of-function recommendations would be considerably strengthened by increasing coverage from 60\% to 90\%+ of cell lines. Set aside two to three months for retraining and data integration.

Priority 3: Stratified cross-validation for tissue-specific reliability evaluation. The mutation landscape varies depending on the tissue. Within-tissue ranking would be enhanced by tissue-specific weights. Set aside one to two months for analysis.

Priority 4: Extension to non-cancer situations (primary cells, organoids, PDX) is the fourth priority. This significantly increases the use case. Give data gathering and validation four to six months.

Priority 5: Using larger validation cohorts to improve default weights. Weights should only be improved if Priority 1 validation is successful. After the experimental validation is finished, set aside two to three months.

\section{Team Contribution and Collaboration Dynamics}

Four team members from complementary specialties worked together on the project, with a defined division of labor and close integration.

{\bf Karen Coutinho (Data Integration and Analysis).} Oversaw the resolution of nomenclature and oversaw the harmonization of five public data sources onto canonical identifiers (DepMap ach\_id, Cellosaurus CVCL, gene symbols, Ensembl IDs). Three significant data handling errors (HPA/GEO identification misattribution, non-deterministic tie-breaking and score-saturation artifact) were found and fixed. authored the final dissertation and oversaw the validation strategy. This position requires a thorough understanding of validation methods, cleaning logic and data provenance.

{\bf Ishaan Bhalla (Graph-Based Ranking).} Neo4j knowledge graph was created with 2,076 cell lines 255 KEGG pathways and 19,000 genes. For graph‑based ranking Neo4j knowledge graph used a walk with restart that incorporated a customized PageRank. The study examined pathway‑coherence score, with four techniques and the results were recorded in a negative‑result section.The author contributed to the design of the architecture of the Neo4j knowledge graph.This study showed why transcriptomic pathway approaches are unable to identify activation and how graph topologies of the Neo4j knowledge graph might capture biological interactions.

{\bf Aman Raj (System Architecture and API).} FastAPI backend was built with four routers: genes, ranking, cell lines and filters. FastAPI backend was designed for access. Data pipeline and Parquet was implemented to perform queries on 222 million rows of data. CI/CD setup was done with GitHub Actions ensuring that solution was not a research project but also a production-ready application.

{\bf Dixit Kaloorani Malarmannan (Web Frontend).} Search Bar, Filter Panel, Results Table and Detail Panel components were added to a React frontend. Visualizing evidence for each source and being able to filter in time was implemented. A PDF export was created to prepare reports, for publication. This work made the system easy for non-technical users to understand and created a visual display of the evidence that domain specialists could read and understand.

{\bf Dr. Daniel D'Andrea (Supervision).} Provided advice on validation tactics, experimental design and mechanistic interpretation of findings. It was especially useful when discussing the limits of techniques, for detecting post‑translational activation the ways pathway investigations can fail and gene‑class hypotheses. He oversaw the process focusing closely on biological accuracy, reproducibility and rigor.

{\bf Collaboration With AstraZeneca.} Two important contributions were made by the industry partner. First problem motivation based on research processes. Cell-line selection is a bottleneck that the system resolves. Second independent verification that revealed flaws. The TP53 wild-type penalty problem was discovered by the domain expert testing the system, which would not have been discovered by measurements alone. This result highlights how important industrial collaboration is to research. The study topic and the body of evidence are strengthened by real-world limitations, which also generate depth and practical validation. The feedback loop was quick: the team. Made corrections within weeks after the researchers presented their findings and a domain expert found issues within days.

{\bf Knowledge Synthesis.} We observed that Bioinformatics (harmonization of -omics data) data engineering (Parquet, ETL pipelines) machine learning (ranking, optimization cross validation) software engineering (API design, CI/CD, deployment) biology (classes of genes pathways, tumor suppressors) and domain expertise (workflows for cell line selection) were all essential for the project. No team member possessed expertise in Bioinformatics, data engineering, machine learning, software engineering, biology and domain expertise. Frequent communication and information exchange were crucial for success. Weekly meetings were not devoted to status reporting; instead we focused on discussing issues and decisions. Asynchronous communication about issues, via GitHub enabled us to have technical discussions while keeping all the context intact.

\section{Closing Reflections and Path Forward}

CellSelector Omics demonstrates that combining data from sources can produce very useful scientific tools when the work is done carefully and checked against expert knowledge. CellSelector Omics’ main contribution is the discovery that gene class, not source-level weighting's the key factor that determines ranking accuracy. This insight came from a study of cases that failed and it reveals a mechanism that goes deeper than the technology itself. Regardless of the method used CellSelector Omics shows that it is important to know why different genes need different levels of evidence.

The final system was shaped directly by CellSelector Omics’ commitment to interpretability. Every recommendation can be traced back to sources of evidence. Every ranking decision can be split into named parts. Adoption of CellSelector Omics in research settings and validation by expert reviewers rely on this transparency.

The study of outcomes—failed pathway‑coherence scoring methods—is a clear example of scientific rigor. Than abandoning the strategy without an explanation CellSelector Omics investigated the reasons why all four methods failed and connected those reasons to simple features of post‑translational activation mechanisms. This understanding guided the choice of an alternative, mutation‑impact scoring and will help future researchers who face similar problems.

A problem with using only numbers was shown when expert reviewers found the TP53 wild‑type penalty flaw. CellSelector Omics can check consistency with metrics such as MRR and cross‑validation. However these metrics cannot prove accuracy; studies of tumor suppressors should use wild‑type TP53. This lesson applies to all bioinformatics tools: do not rely on experts at the end; involve them early and often.

The partnership with AstraZeneca proved useful. The problem arose from research limits. The solution was confirmed by domain experts. Quick changes were possible because of feedback loops. When such collaborations are set up correctly academic research gains a lot: the company gets a tool the scientists get real‑world grounding and confirmation and science moves faster.

The best next steps are clear. CellSelector Omics’ recommendations will be tested in situations through experiments. Loss‑of‑function rankings will improve when more mutation data is added. CellSelector Omics will also be useful, in non‑cancer settings. With the people and resources these goals can be achieved in twelve to eighteen months.

% =============================================================================

\backmatter

\bibliography{references}

\chapter*{Appendix}
\addcontentsline{toc}{chapter}{Appendix}

The deployed platform is publicly accessible at:
\begin{center}
    \url{https://cellselectoromics.com/}
\end{center}

\vspace{-0.2cm}
The full source code is available at:
\begin{center}
    \url{https://github.com/karencoutinhoxx/cellselector-omics}
\end{center}

\vspace{0.3cm}



% =============================================================================

\end{document}